
% BEGIN S1
 
\par\medskip\noindent\textbf{1 Introduction}\par
 
% BEGIN S1.SS0.SSS0.Px1
 
\par\medskip\noindent\textbf{Motivation}\par
 
% BEGIN S1.SS0.SSS0.Px1.p1
 
Quantum computing is an emerging paradigm of computation, where quantum physical phenomena such as entanglement and superposition are used to obtain an advantage over classical computation. A testament to the richness of the field is the variety of computational models: quantum Turing machines [6], quantum circuits [31, 26], measurement-based quantum computation [7, 12], linear optical circuits [21], among others.\par
 
% END S1.SS0.SSS0.Px1.p1
 
% BEGIN S1.SS0.SSS0.Px1.p2
 
In recent decades, a lot of effort has been put into developing high-level quantum programming languages that allow programmers to abstract away from the technicalities of these low-level models [27, 3]. Toward that end, several verification techniques such as type systems [15] or categorical approaches for reasoning about program semantics [4, 19] have been studied and developed to ensure the physical reality of compiled programs, for example, by ensuring that they satisfy the properties of quantum mechanics such as the no-cloning theorem [1] and unitarity [14].\par
 
% END S1.SS0.SSS0.Px1.p2
 
% BEGIN S1.SS0.SSS0.Px1.p3
 
An important line of research in this area involves checking polytime termination of quantum programs [11, 29, 17], by showing that any program can be simulated by a polytime quantum Turing machine (QTM). As a consequence of Yao’s Theorem [31], demonstrating that polynomial-time QTMs are computationally equivalent to uniform and poly-sized quantum circuit families, such programs can be instantiated by a uniform family of quantum circuits of polynomial size, i.e., with a polynomial number of qubits and gates. However, direct compilation techniques often fail to ensure this property, due to the use of quantum branching, where the flow in a conditional is determined by the state of a qubit.\par
 
% END S1.SS0.SSS0.Px1.p3
 
% BEGIN S1.SS0.SSS0.Px1.p4
 
As an illustrative example, consider a quantum case, of the shape\par
    \[\text{{{qcase }}}\mathrm{q}_{1},\dots,\mathrm{q}_{k}\text{{{ of }}}\{i\to%
\mathrm{S}_{i}\}_{i\in\{0,1\}^{k}},\]    
which executes the superposition of the unitary transformations $U_{i}$ implemented by statements $\mathrm{S}_{i}$, depending on the state $\left|i\right\rangle$ of the control qubits $\mathrm{q}_{1},\dots,\mathrm{q}_{k}$. This statement can be simulated by a QTM that, depending on $k$ symbols in its tape, performs the instructions of different QTMs $M_{i}$ in superposition. Consequently, its runtime is the maximum runtime among all $M_{i}$ ([6, Branching Lemma]). Instantiating this qcase statement on quantum circuits can be done sequentially (in-depth) or in parallel (in-width). An in-depth encoding consists in applying each controlled-$U_{i}$ gate sequentially, thus implementing the unitary transformation $\sum_{i\in\{0,1\}^{k}}\outerproduct{i}{i}\otimes U_{i}$. The depth of the resulting circuit is the sum of the depths of all $U_{i}$, which is exponential in the number $k$ of control qubits. An alternative strategy, inspired by QRAM, allows for the execution of each $U_{i}$ in parallel in the circuit, using controlled swaps to perform routing of qubit addresses in linear depth on the number of control qubits. This strategy has been recently applied in compilation techniques to automatically parallelize the branches of qcase statements and allow for optimal parallel time (depth) complexity [35].\par
 
% END S1.SS0.SSS0.Px1.p4
 
% BEGIN S1.SS0.SSS0.Px1.p5
 
Since the unitaries are parallelized, the circuit depth is the maximum among the depths of $U_{i}$, but this strategy results in circuits where the width (number of qubits) scales as the sum of the complexities of $U_{i}$. Consequently, this parallelization requires a number of ancillas that grows exponentially on the number of control qubits. In summary, any implementation, that is agnostic about the structure of each $U_{i}$ will incur such an exponential cost; this is the case of qcase statements used in the preparation of arbitrary quantum states [34, 16].\par
 
% END S1.SS0.SSS0.Px1.p5
 
% BEGIN S1.SS0.SSS0.Px1.p6
 
Automatic implementations of quantum conditionals therefore produce circuits with a size that scales with the sum (rather than the maximum) of the complexity of each branch. This problem, coined branch sequentialization in [32], leads, in many cases, to an exponential blow-up in circuit complexity, forcing the programmer to rewrite and optimize their code in order to make use of the program structure and improve the complexity bound. A challenging issue is, therefore, to develop quantum programming languages that avoid branch sequentialization by ensuring the correct circuit complexity for quantum control statements while providing full use of quantum control to the programmer [32, 33].\par
 
% END S1.SS0.SSS0.Px1.p6
 
% END S1.SS0.SSS0.Px1
 
% BEGIN S1.SS0.SSS0.Px2
 
\par\medskip\noindent\textbf{Contribution}\par
 
% BEGIN S1.SS0.SSS0.Px2.p1
 
This paper solves the above issue by introducing a programming language, called pbp for Polynomially Branching Program, with quantum case and first-order recursive procedures, on which compilation of quantum conditional does not lead to an exponential blow-up in circuit size. Our main contributions are as follows:\par
  
\par\noindent  $\blacksquare$ 
% BEGIN S1.I1.i1.p1
 
We introduce a compilation strategy compile (Figures 8 and 9) into quantum circuits, show its soundness (Theorem 8), and show that it avoids branch sequentialization on pbp (Theorem 9), a language with restrictions on recursive calls to procedures. Toward that end, we formally define the time complexity $\textnormal{Time}_{\mathrm{P}}:\mathbb{N}\to\mathbb{N}$ of a program $\mathrm{P}$ (Definition 1) as a map from the number $n$ of qubits to the maximal number of procedures called during a program execution on the Hilbert space $\mathbb{C}^{2^{n}}$. The time complexity of a qcase statement is precisely the maximum of the time complexities of its branches and we show that branch sequentialization is avoided as the circuits generated by compile have size bounded asymptotically by the time complexity of the compiled program (Theorem 9).\par
 
% END S1.I1.i1.p1
 
 
\par\noindent  $\blacksquare$ 
% BEGIN S1.I1.i2.p1
 
A natural question concerns the expressive power of language pbp. We also solve this issue by showing that pbp is sound and complete for quantum polynomial time (Theorem 11). Consequently, any polytime quantum algorithm can be encoded by a pbp program. We illustrate the methodology through well-known examples such as the program QFT realizing the quantum Fourier transform (Example 6).\par
 
% END S1.I1.i2.p1
 
 
\par\noindent  $\blacksquare$ 
% BEGIN S1.I1.i3.p1
 
We also discuss asymptotic bounds for a strict extension of pbp, where the reduction in the input sorted set passed to procedure calls can be arbitrary. We show that this extension only increases the circuit complexity of the no-branch-sequentialization case by a linear factor (Theorem 13).\par
 
% END S1.I1.i3.p1
 
 
\par\noindent  $\blacksquare$ 
% BEGIN S1.I1.i4.p1
 
We also show that compile strictly improves on existing compilation strategies on first-order quantum programs with quantum case [17] by a polynomial speedup of arbitrarily large degree (Table 1).\par
 
% END S1.I1.i4.p1
 
  
% END S1.SS0.SSS0.Px2.p1
 
% END S1.SS0.SSS0.Px2
 
% BEGIN S1.SS0.SSS0.Px3
 
\par\medskip\noindent\textbf{A bird’s-eye view of our compilation strategy}\par
 
% BEGIN S1.SS0.SSS0.Px3.p1
 
We will now consider a simple illustrating example where branch sequentialization can exponentially worsen the size complexity of the compiled circuit, and show how our compilation strategy can preserve a polynomial bound.\par
 
% END S1.SS0.SSS0.Px3.p1
 
% BEGIN S1.SS0.SSS0.Px3.p2
 
The program PAIRS defined in Figure 1 consists in a simple call to procedure pairs (line 7) on a list $\bar{\mathrm{q}}$ of unique qubits. Let $|\bar{\mathrm{q}}|$ be the number of qubits in $\bar{\mathrm{q}}$. By language design, the procedure pairs immediately terminates whenever $|\bar{\mathrm{q}}|=0$. The procedure enters a quantum case (line 2) when $|\bar{\mathrm{q}}|\geq 2$, otherwise it applies a NOT gate to the single qubit (line 6). On line 2, the program will branch depending on the state $\bar{\mathrm{q}}\textnormal{\small[$1,2$]}$ of the first two qubits in $\bar{\mathrm{q}}$. Out of all four cases (lines 2-5), pairs only performs an operation when the first two qubits are in state $\left|00\right\rangle$ or $\left|11\right\rangle$, in which case it performs a recursive call on $\bar{\mathrm{q}}\ominus\textnormal{\small[$1,2$]}$, the list obtained by removing the first and second qubits of $\bar{\mathrm{q}}$.\par
 
% END S1.SS0.SSS0.Px3.p2
 
% BEGIN S1.F1

\begin{center}\includegraphics[width=0.98\linewidth,height=0.70\textheight,keepaspectratio]{graphic-01.pdf}\end{center}
 
Figure 1: Program PAIRS.\par
 
% END S1.F1
 
% BEGIN S1.SS0.SSS0.Px3.p3
 
Given an input state $\left|xy\right\rangle$, with $x\in\{0,1\}^{*}$ and $y\in\{0,1\}$, pairs will apply a NOT gate to the last qubit $y$ if and only if $x$ is a string in the language $(00\mid 11)^{\ast}$. Since pairs performs at most one call per branch, and consumes 2 qubits from its input while doing so, its time complexity is in $O(|\bar{\mathrm{q}}|)$, when taking the qcase complexity as the maximum of each branch.\par
 
% END S1.SS0.SSS0.Px3.p3
 
% BEGIN S1.SS0.SSS0.Px3.p4
 
We now discuss the circuit compilation of pairs, when $|\bar{\mathrm{q}}|\geq 2$. In Figure 2, we give three strategies (a), (b), and (c), exemplifying different approaches. Strategies (a) and (b) are automatic compilation strategies, which ignore the inner structure of the program – the fact that the bodies of the two non-skip branches of the qcase are identical and therefore encode the same unitary transformation. Figure 2(a) represents an in-depth strategy, whereas Figure 2(b) implements an in-width strategy, making use of different registers $r_{00}$ and $r_{11}$, both initialized to $\left|0^{|\bar{\mathrm{q}}|-2}\right\rangle$ to perform each branch in parallel and then recombine the results in the same register. These two implementations require two recursive calls to pairs and, consequently, their size complexity is the sum of the complexities of the two branches. These strategies are performing branch sequentialization and, in both cases, the compiled circuit has size exponential in $|\bar{\mathrm{q}}|$. In contrast, Figure 2(c) avoids this exponential blow-up, making use of the fact that the two branches are identical. This strategy is able to merge the recursive calls into a single procedure call with the use of one ancilla. This is precisely the type of strategy used by our compilation algorithm compile (Section 3). Although this example is relatively simple due to similar recursive calls, compile allows us to deal with more involved situations, where branches are not identical.\par
 
% END S1.SS0.SSS0.Px3.p4
 
% BEGIN S1.F2

\begin{center}\includegraphics[width=0.98\linewidth,height=0.70\textheight,keepaspectratio]{graphic-02.pdf}\end{center}
 
Figure 2: Compilation strategies.\par
 
% END S1.F2
 
% END S1.SS0.SSS0.Px3
 
% BEGIN S1.SS0.SSS0.Px4
 
\par\medskip\noindent\textbf{Related work}\par
 
% BEGIN S1.SS0.SSS0.Px4.p1
 
Resource optimization in low-level models of quantum computation is a well-studied subject. Given a specific quantum circuit, it is possible to reduce its number of gates (or at least its number of non-Clifford gates) via techniques such as gate substitution and graph rewriting [25, 23, 20, 13]. The study of resource optimization in the asymptotic scenario is a relatively young research area as it involves reasoning about families of circuits and program structure. Different implicit characterizations of quantum complexity classes have been developed using a lambda-calculus [11], function algebras [29, 30] and a first-order programming language [17]. The last of these provided a compilation strategy into quantum circuits with bounds on the circuit size based on the syntactic restrictions placed on the programs.\par
 
% END S1.SS0.SSS0.Px4.p1
 
% BEGIN S1.SS0.SSS0.Px4.p2
 
Compilation strategies for the quantum control statement (also called quantum switch case or quantum multiplexer) have been studied, for instance, in state preparation [22], appearing in quantum machine learning [18] and Hamiltonian simulation [8]. These techniques either focus on optimizing number of qubits (circuit width) [34] or circuit depth [28], but in all cases correspond to circuits of exponential size. In order to improve on these bounds, one must restrict the set of programs to those that admit an efficient implementation, which can be deduced from the program structure. Optimized compilation techniques in that scenario can then be judged on expressivity and completeness: how easily can one write a program while ensuring the syntactical restrictions? Do these restrictions encompass all efficient programs? The field of implicit computational complexity is particularly well-suited to answer these questions [9, 10].\par
 
% END S1.SS0.SSS0.Px4.p2
 
% END S1.SS0.SSS0.Px4
 
% END S1


% BEGIN S2
 
\par\medskip\noindent\textbf{2 First-Order Programming Language with Quantum Case}\par
 
% BEGIN S2.p1
 
In this section, we introduce a programming language with quantum control and first-order recursive procedures. After introducing its syntax and semantics, we introduce a restriction pbp (for Polynomially Branching Programs) on which branch sequentialization can be avoided.\par
 
% END S2.p1
 
% BEGIN S2.F3
 
$\begin{array}[]{lllll}\text{(Integers)}&\mathrm{i},\mathrm{j},\mathrm{k}&::=&%
\mathrm{x}\ |\ n\ |\ \mathrm{i}\pm n\ |\ |\mathrm{s}|\\
\text{(Booleans)}&\mathrm{b}&::=&\mathrm{i}\geq\mathrm{i}\ |\ \cdots\ |\ %
\mathrm{b}\wedge\mathrm{b}\ |\ \cdots\\
\text{(Qubits)}&\mathrm{q}&::=&\mathrm{s}\textnormal{\small[$\mathrm{i}$]}\\
\text{(Sorted sets)}&\mathrm{s}&::=&\bar{\mathrm{q}}\ |\ \mathrm{s}\ominus%
\textnormal{\small[$\mathrm{i}$]}\\
\text{(Statements)}&\mathrm{S}&::=&\text{\color[rgb]{%
0.03515625,0.0390625,0.72265625}\definecolor[named]{pgfstrokecolor}{rgb}{%
0.03515625,0.0390625,0.72265625}{{ skip}}};\ |\ \mathrm{q}{\ \textnormal{{%
\textasteriskcentered=}}\ }\mathrm{U}^{f}(\mathrm{j});\ |\ \mathrm{S}\ \mathrm%
{S}\ |\ \text{\color[rgb]{0.03515625,0.0390625,0.72265625}\definecolor[named]{%
pgfstrokecolor}{rgb}{0.03515625,0.0390625,0.72265625}{{if }}}\mathrm{b}\text{%
\color[rgb]{0.03515625,0.0390625,0.72265625}\definecolor[named]{pgfstrokecolor%
}{rgb}{0.03515625,0.0390625,0.72265625}{{ then }}}\mathrm{S}\text{\color[rgb]{%
0.03515625,0.0390625,0.72265625}\definecolor[named]{pgfstrokecolor}{rgb}{%
0.03515625,0.0390625,0.72265625}{{ else }}}\mathrm{S}\\
&&&\ |\ \text{\color[rgb]{0.03515625,0.0390625,0.72265625}\definecolor[named]{%
pgfstrokecolor}{rgb}{0.03515625,0.0390625,0.72265625}{{qcase }}}\mathrm{q}%
\text{\color[rgb]{0.03515625,0.0390625,0.72265625}\definecolor[named]{%
pgfstrokecolor}{rgb}{0.03515625,0.0390625,0.72265625}{{ of }}}\{0\to\mathrm{S}%
,1\to\mathrm{S}\}\ |\ \text{\color[rgb]{0.03515625,0.0390625,0.72265625}%
\definecolor[named]{pgfstrokecolor}{rgb}{0.03515625,0.0390625,0.72265625}{{%
call }}}\text{\color[rgb]{0.79296875,0.3203125,0.8359375}\definecolor[named]{%
pgfstrokecolor}{rgb}{0.79296875,0.3203125,0.8359375}{{proc}}}(\mathrm{s});&\\
\text{(Procedure declarations)}&\mathrm{D}&::=&\ \varepsilon\ |\ \text{\color[%
rgb]{0.03515625,0.0390625,0.72265625}\definecolor[named]{pgfstrokecolor}{rgb}{%
0.03515625,0.0390625,0.72265625}{{decl }}}\text{\color[rgb]{%
0.79296875,0.3203125,0.8359375}\definecolor[named]{pgfstrokecolor}{rgb}{%
0.79296875,0.3203125,0.8359375}{{proc}}}(\bar{\mathrm{q}})\{\mathrm{S}\},\,%
\mathrm{D}\\
\text{(Programs)}&\mathrm{P}(\bar{\mathrm{q}})&::=&\,\mathrm{D}::\mathrm{S}%
\end{array}$\par
 
Figure 3: Syntax of programs.\par
 
% END S2.F3
 
% BEGIN S2.SS1
 
\par\medskip\noindent\textbf{2.1 Syntax of Programs}\par
 
% BEGIN S2.SS1.p1
 
The language includes four basic datatypes for expressions, whose corresponding expressions are described in Figure 3.\par
  
\par\noindent  $\blacksquare$ 
% BEGIN S2.I1.i1.p1
 
Integers: Integer expressions are variables $\mathrm{x}$, constant $n\in\mathbb{N}$, an addition by a constant $\mathrm{i}\pm n$, or the size $|\mathrm{s}|$ of a sorted set $\mathrm{s}$.\par
 
% END S2.I1.i1.p1
 
 
\par\noindent  $\blacksquare$ 
% BEGIN S2.I1.i2.p1
 
Booleans: Such expressions $\mathrm{b}$ are defined in a standard way.\par
 
% END S2.I1.i2.p1
 
 
\par\noindent  $\blacksquare$ 
% BEGIN S2.I1.i3.p1
 
Qubit: qubits expressions are of the shape $\mathrm{s}\textnormal{\small[$\mathrm{i}$]}$ which denotes the $\mathrm{i}$-th qubit in sorted set $\mathrm{s}$.\par
 
% END S2.I1.i3.p1
 
 
\par\noindent  $\blacksquare$ 
% BEGIN S2.I1.i4.p1
 
Sorted sets: lists of unique (i.e., non-duplicable) qubit pointers. Sorted-set expressions $\mathrm{s}$ are either variables $\bar{\mathrm{q}}$ or $\mathrm{s}\ominus\textnormal{\small[$\mathrm{i}$]}$, the sorted set $\mathrm{s}$ where the $\mathrm{i}$-th element has been removed.\par
 
% END S2.I1.i4.p1
 
  
Let $\mathrm{s}\textnormal{\small[$\mathrm{i}_{1},\ldots,\mathrm{i}_{n}$]}$ be a shorthand for $\mathrm{s}\textnormal{\small[$\mathrm{i}_{1}$]},\ldots,\mathrm{s}\textnormal{%
\small[$\mathrm{i}_{n}$]}$. We also define syntactic sugar for pointing to the $n$-th last qubit in a sorted set, by defining for any $n\geq 1$, $\bar{\mathrm{q}}\textnormal{\small[$-n$]}\triangleq\bar{\mathrm{q}}\textnormal%
{\small[$|\bar{\mathrm{q}}|-n+1$]}$.\par
 
% END S2.SS1.p1
 
% BEGIN S2.SS1.p2
 
A program $\mathrm{P}\triangleq\mathrm{D}::\mathrm{S}$ is defined in Figure 3 as a list of (possibly recursive) procedure declarations $\mathrm{D}$, followed by a program statement $\mathrm{S}$.\par
 
% END S2.SS1.p2
 
% BEGIN S2.SS1.p3
 
Let Procedures be an enumerable set of procedure names. We write $\text{\color[rgb]{0.79296875,0.3203125,0.8359375}\definecolor[named]{%
pgfstrokecolor}{rgb}{0.79296875,0.3203125,0.8359375}{{proc}}}\in\mathrm{P}$ to denote that proc appears in $\mathrm{P}$. Each procedure of name $\text{\color[rgb]{0.79296875,0.3203125,0.8359375}\definecolor[named]{%
pgfstrokecolor}{rgb}{0.79296875,0.3203125,0.8359375}{{proc}}}\in\mathrm{P}$ is defined by a (unique) procedure declaration $\text{\color[rgb]{0.03515625,0.0390625,0.72265625}\definecolor[named]{%
pgfstrokecolor}{rgb}{0.03515625,0.0390625,0.72265625}{{decl }}}\text{\color[%
rgb]{0.79296875,0.3203125,0.8359375}\definecolor[named]{pgfstrokecolor}{rgb}{%
0.79296875,0.3203125,0.8359375}{{proc}}}(\bar{\mathrm{q}})\{\mathrm{S}\}\in%
\mathrm{D}$, which takes a sorted set $\bar{\mathrm{q}}$ as input parameter and executes the procedure statement $\mathrm{S}$. We sometimes write $\mathrm{S}^{\text{\color[rgb]{0.79296875,0.3203125,0.8359375}\definecolor[%
named]{pgfstrokecolor}{rgb}{0.79296875,0.3203125,0.8359375}{{proc}}}}$ to explicitly refer to the procedure statement $\mathrm{S}$ of proc.\par
 
% END S2.SS1.p3
 
% BEGIN S2.SS1.p4
 
Statements include the no-op instruction, unitary application, sequences, conditional, quantum case, and procedures calls. For the sake of universality [5], in a unitary application $\mathrm{q}{\ \textnormal{{\textasteriskcentered=}}\ }\mathrm{U}^{f}(\mathrm{j});$, $\mathrm{U}^{f}(\mathrm{j})$ is a unitary transformation that can take an integer $\mathrm{j}$ and a polynomial-time approximable [2] total function $f\in\mathbb{Z}\to[0,2\pi)$ as optional arguments. The $f$ and $\mathrm{i}$ can be omitted when they are not useful, as in a NOT gate.\par
 
% END S2.SS1.p4
 
% BEGIN S2.SS1.p5
 
The quantum conditional $\text{\color[rgb]{0.03515625,0.0390625,0.72265625}\definecolor[named]{%
pgfstrokecolor}{rgb}{0.03515625,0.0390625,0.72265625}{{qcase }}}\mathrm{s}%
\textnormal{\small[$\mathrm{i}$]}\text{\color[rgb]{%
0.03515625,0.0390625,0.72265625}\definecolor[named]{pgfstrokecolor}{rgb}{%
0.03515625,0.0390625,0.72265625}{{ of }}}\{0\to\mathrm{S}_{0},1\to\mathrm{S}_{%
1}\}$ allows branching by executing statements $\mathrm{S}_{0}$ and $\mathrm{S}_{1}$ in superposition according to the state of qubit $\mathrm{s}[\mathrm{i}]$. The procedure call $\text{\color[rgb]{0.03515625,0.0390625,0.72265625}\definecolor[named]{%
pgfstrokecolor}{rgb}{0.03515625,0.0390625,0.72265625}{{call }}}\text{\color[%
rgb]{0.79296875,0.3203125,0.8359375}\definecolor[named]{pgfstrokecolor}{rgb}{%
0.79296875,0.3203125,0.8359375}{{proc}}}(\mathrm{s});$ runs procedure proc with sorted set expression $\mathrm{s}$. The quantum conditional can be extended to multiple qubits in a standard way as used in Figure 1: $\text{\color[rgb]{0.03515625,0.0390625,0.72265625}\definecolor[named]{%
pgfstrokecolor}{rgb}{0.03515625,0.0390625,0.72265625}{{qcase }}}\mathrm{q}_{1}%
,\mathrm{q}_{2}\text{\color[rgb]{0.03515625,0.0390625,0.72265625}\definecolor[%
named]{pgfstrokecolor}{rgb}{0.03515625,0.0390625,0.72265625}{{ of }}}\{00\to%
\mathrm{S}_{00},01\to\mathrm{S}_{01},10\to\mathrm{S}_{10},11\to\mathrm{S}_{11}\}$ is a shorthand notation for $\text{\color[rgb]{0.03515625,0.0390625,0.72265625}\definecolor[named]{%
pgfstrokecolor}{rgb}{0.03515625,0.0390625,0.72265625}{{qcase }}}\mathrm{q}_{1}%
\text{\color[rgb]{0.03515625,0.0390625,0.72265625}\definecolor[named]{%
pgfstrokecolor}{rgb}{0.03515625,0.0390625,0.72265625}{{ of }}}\{0\to\text{%
\color[rgb]{0.03515625,0.0390625,0.72265625}\definecolor[named]{pgfstrokecolor%
}{rgb}{0.03515625,0.0390625,0.72265625}{{qcase }}}\mathrm{q}_{2}\text{\color[%
rgb]{0.03515625,0.0390625,0.72265625}\definecolor[named]{pgfstrokecolor}{rgb}{%
0.03515625,0.0390625,0.72265625}{{ of }}}\{0\to\mathrm{S}_{00},1\to\mathrm{S}_%
{01}\},1\to\text{\color[rgb]{0.03515625,0.0390625,0.72265625}\definecolor[%
named]{pgfstrokecolor}{rgb}{0.03515625,0.0390625,0.72265625}{{qcase }}}\mathrm%
{q}_{2}\text{\color[rgb]{0.03515625,0.0390625,0.72265625}\definecolor[named]{%
pgfstrokecolor}{rgb}{0.03515625,0.0390625,0.72265625}{{ of }}}\{0\to\mathrm{S}%
_{10},1\to\mathrm{S}_{11}\}\}$.\par
 
% END S2.SS1.p5
 
% BEGIN S2.SS1.p6
 
We also make use of some syntactic sugar to describe statements encoding constant-time quantum operations. For instance, the CNOT, SWAP, as well as a controlled-phase shift gate are defined by:\par
    $\displaystyle\text{CNOT}(\mathrm{q}_{1},\mathrm{q}_{2})\ \triangleq$ $\displaystyle\ \text{\color[rgb]{0.03515625,0.0390625,0.72265625}\definecolor[%
named]{pgfstrokecolor}{rgb}{0.03515625,0.0390625,0.72265625}{{qcase }}}\mathrm%
{q}_{1}\text{\color[rgb]{0.03515625,0.0390625,0.72265625}\definecolor[named]{%
pgfstrokecolor}{rgb}{0.03515625,0.0390625,0.72265625}{{ of }}}\{0\to\text{%
\color[rgb]{0.03515625,0.0390625,0.72265625}\definecolor[named]{pgfstrokecolor%
}{rgb}{0.03515625,0.0390625,0.72265625}{{skip}}};,\,1\to\mathrm{q}_{2}{\ %
\textnormal{{\textasteriskcentered=}}\ }\text{NOT};\}$     $\displaystyle\text{SWAP}(\mathrm{q}_{1},\mathrm{q}_{2})\ \triangleq$ $\displaystyle\ \text{CNOT}(\mathrm{q}_{1},\mathrm{q}_{2})\ \text{CNOT}(\mathrm%
{q}_{2},\mathrm{q}_{1})\ \text{CNOT}(\mathrm{q}_{1},\mathrm{q}_{2})$     $\displaystyle\textnormal{CPHASE}(\mathrm{q}_{1},\,\mathrm{q}_{2},\,\mathrm{i})\triangleq$ $\displaystyle\text{\color[rgb]{0.03515625,0.0390625,0.72265625}\definecolor[%
named]{pgfstrokecolor}{rgb}{0.03515625,0.0390625,0.72265625}{{ qcase }}}%
\mathrm{q}_{1}\text{\color[rgb]{0.03515625,0.0390625,0.72265625}\definecolor[%
named]{pgfstrokecolor}{rgb}{0.03515625,0.0390625,0.72265625}{{ of }}}\{0\to%
\text{\color[rgb]{0.03515625,0.0390625,0.72265625}\definecolor[named]{%
pgfstrokecolor}{rgb}{0.03515625,0.0390625,0.72265625}{{skip}}};,1\to\mathrm{q}%
_{2}{\ \textnormal{{\textasteriskcentered=}}\ }\text{Ph}^{\lambda x.\pi/2^{x-1%
}}(\mathrm{i});\}$    
% END S2.SS1.p6
 
% BEGIN S2.SS1.p7
 
Given a program $\mathrm{P}\triangleq\mathrm{D}::\mathrm{S}$, the call relation $\to_{\mathrm{P}}\ \subseteq\textnormal{Procedures}\times\textnormal{Procedures}$ is defined for any two procedures $\text{\color[rgb]{0.79296875,0.3203125,0.8359375}\definecolor[named]{%
pgfstrokecolor}{rgb}{0.79296875,0.3203125,0.8359375}{{proc${}_{1}$}}},$ $\text{\color[rgb]{0.79296875,0.3203125,0.8359375}\definecolor[named]{%
pgfstrokecolor}{rgb}{0.79296875,0.3203125,0.8359375}{{proc${}_{2}$}}}\in%
\mathrm{P}$ as $\text{\color[rgb]{0.79296875,0.3203125,0.8359375}\definecolor[named]{%
pgfstrokecolor}{rgb}{0.79296875,0.3203125,0.8359375}{{proc${}_{1}$}}}\to_{%
\mathrm{P}}\text{\color[rgb]{0.79296875,0.3203125,0.8359375}\definecolor[named%
]{pgfstrokecolor}{rgb}{0.79296875,0.3203125,0.8359375}{{proc${}_{2}$}}}$ whenever $\text{\color[rgb]{0.79296875,0.3203125,0.8359375}\definecolor[named]{%
pgfstrokecolor}{rgb}{0.79296875,0.3203125,0.8359375}{{proc${}_{2}$}}}\in%
\mathrm{S}^{\text{\color[rgb]{0.79296875,0.3203125,0.8359375}\definecolor[%
named]{pgfstrokecolor}{rgb}{0.79296875,0.3203125,0.8359375}{{proc${}_{1}$}}}}$. The relation $\succeq_{\mathrm{P}}$ is then the transitive closure of $\to_{\mathrm{P}}$, and $\sim_{\mathrm{P}}$ denotes the equivalence relation where $\text{\color[rgb]{0.79296875,0.3203125,0.8359375}\definecolor[named]{%
pgfstrokecolor}{rgb}{0.79296875,0.3203125,0.8359375}{{proc${}_{1}$}}}\sim_{%
\mathrm{P}}\text{\color[rgb]{0.79296875,0.3203125,0.8359375}\definecolor[named%
]{pgfstrokecolor}{rgb}{0.79296875,0.3203125,0.8359375}{{proc${}_{2}$}}}$ if $\text{\color[rgb]{0.79296875,0.3203125,0.8359375}\definecolor[named]{%
pgfstrokecolor}{rgb}{0.79296875,0.3203125,0.8359375}{{proc${}_{1}$}}}\succeq_{%
\mathrm{P}}\text{\color[rgb]{0.79296875,0.3203125,0.8359375}\definecolor[named%
]{pgfstrokecolor}{rgb}{0.79296875,0.3203125,0.8359375}{{proc${}_{2}$}}}$ and $\text{\color[rgb]{0.79296875,0.3203125,0.8359375}\definecolor[named]{%
pgfstrokecolor}{rgb}{0.79296875,0.3203125,0.8359375}{{proc${}_{2}$}}}\succeq_{%
\mathrm{P}}\text{\color[rgb]{0.79296875,0.3203125,0.8359375}\definecolor[named%
]{pgfstrokecolor}{rgb}{0.79296875,0.3203125,0.8359375}{{proc${}_{1}$}}}$ both hold.\par
 
% END S2.SS1.p7
 
% BEGIN S2.SS1.p8
 
A procedure proc is recursive whenever $\text{\color[rgb]{0.79296875,0.3203125,0.8359375}\definecolor[named]{%
pgfstrokecolor}{rgb}{0.79296875,0.3203125,0.8359375}{{proc}}}\sim_{\mathrm{P}}%
\text{\color[rgb]{0.79296875,0.3203125,0.8359375}\definecolor[named]{%
pgfstrokecolor}{rgb}{0.79296875,0.3203125,0.8359375}{{proc}}}$ holds. The strict order $\succ_{\mathrm{P}}$ is defined as $\text{\color[rgb]{0.79296875,0.3203125,0.8359375}\definecolor[named]{%
pgfstrokecolor}{rgb}{0.79296875,0.3203125,0.8359375}{{proc${}_{1}$}}}\succ_{%
\mathrm{P}}\text{\color[rgb]{0.79296875,0.3203125,0.8359375}\definecolor[named%
]{pgfstrokecolor}{rgb}{0.79296875,0.3203125,0.8359375}{{proc${}_{2}$}}}$ if $\text{\color[rgb]{0.79296875,0.3203125,0.8359375}\definecolor[named]{%
pgfstrokecolor}{rgb}{0.79296875,0.3203125,0.8359375}{{proc${}_{1}$}}}\succeq_{%
\mathrm{P}}\text{\color[rgb]{0.79296875,0.3203125,0.8359375}\definecolor[named%
]{pgfstrokecolor}{rgb}{0.79296875,0.3203125,0.8359375}{{proc${}_{2}$}}}$ and $\text{\color[rgb]{0.79296875,0.3203125,0.8359375}\definecolor[named]{%
pgfstrokecolor}{rgb}{0.79296875,0.3203125,0.8359375}{{proc${}_{1}$}}}\not\sim_%
{\mathrm{P}}\text{\color[rgb]{0.79296875,0.3203125,0.8359375}\definecolor[%
named]{pgfstrokecolor}{rgb}{0.79296875,0.3203125,0.8359375}{{proc${}_{2}$}}}$ both hold.\par
 
% END S2.SS1.p8
 
% END S2.SS1
 
% BEGIN S2.SS2
 
\par\medskip\noindent\textbf{2.2 Semantics of Programs}\par
 
% BEGIN S2.SS2.p1
 
Let $\mathbb{B}$ denote the set of Booleans and $\mathcal{L}(\mathbb{N})$ denote the set of lists of natural numbers, $[]$ being the empty list. We interpret each basic type as follows:\par
    $\displaystyle\llbracket\text{Integers}\rrbracket$ $\displaystyle\triangleq\mathbb{Z}$ $\displaystyle\llbracket\text{Booleans}\rrbracket$ $\displaystyle\triangleq\mathbb{B}$ $\displaystyle\llbracket\text{Qubits}\rrbracket$ $\displaystyle\triangleq\mathbb{N}$ $\displaystyle\llbracket\text{Sorted Sets}\rrbracket$ $\displaystyle\triangleq\mathcal{L}(\mathbb{N})$    
Qubits are interpreted as integers (pointers) and sorted sets are interpreted as lists of pointers. For each basic type $\tau$, the semantics of expressions, given in Figure 5, is described in a standard way as a function\par
    \[\Downarrow_{\llbracket\tau\rrbracket}\ \in\tau\times\mathcal{L}(\mathbb{N})\to%
\llbracket\tau\rrbracket.\]    
% END S2.SS2.p1
 
% BEGIN S2.SS2.p2
 
$(\mathrm{e},l)\Downarrow_{\llbracket\tau\rrbracket}v$ holds when expression $\mathrm{e}$ of type $\tau$ evaluates to the value $v\in\llbracket\tau\rrbracket$ under the context $l\in\mathcal{L}(\mathbb{N})$. $l$ is simply the sorted set of qubit pointers considered when evaluating $\mathrm{e}$. For instance, we have that $(\bar{\mathrm{q}}\textnormal{\small[$2$]},[1,4,5])\Downarrow_{\mathbb{N}}4$ (the second qubit has index $4$), $(\bar{\mathrm{q}}\ominus\textnormal{\small[$4$]},[1,4,5])\Downarrow_{\mathcal{%
L}(\mathbb{N})}[]$ ($[]$ is used for errors on type $\mathcal{L}(\mathbb{N})$), $(\bar{\mathrm{q}}\textnormal{\small[$4$]},[1,4,5])\Downarrow_{\mathbb{N}}0$ (index $0$ encodes error on type $\mathbb{N}$), and $(\bar{\mathrm{q}}\ominus\textnormal{\small[$3$]},[1,4,5])\Downarrow_{\mathcal{%
L}(\mathbb{N})}[1,4]$ (third qubit removed).\par
 
% END S2.SS2.p2
 
% BEGIN S2.F4

\begin{center}\includegraphics[width=0.98\linewidth,height=0.70\textheight,keepaspectratio]{graphic-03.pdf}\end{center}
 
Figure 4: Semantics of statements.\par
 
% END S2.F4
 
% BEGIN S2.F5

\begin{center}\includegraphics[width=0.98\linewidth,height=0.70\textheight,keepaspectratio]{graphic-04.pdf}\end{center}
 
Figure 5: Semantics of expressions.\par
 
% END S2.F5
 
% BEGIN S2.SS2.p3
 
Let $\mathcal{H}_{2^{n}}$ denote the Hilbert space of $n$ qubits $\mathbb{C}^{2^{n}}$ and $\mathcal{P}(\mathbb{N})$ denote the powerset of $\mathbb{N}$.\par
 
% END S2.SS2.p3
 
% BEGIN S2.SS2.p4
 
A configuration $c$ over $n$ qubits is of the shape $(\mathrm{S},\left|\psi\right\rangle,A,l)$, for some statement $\mathrm{S}\in\text{Statements}\cup\{\top,\bot\})$, $\top$ and $\bot$ being two special symbols denoting termination and error, respectively, $\left|\psi\right\rangle$ is a quantum state in $\mathcal{H}_{2^{n}}$, where $A\in\mathcal{P}(\mathbb{N})$ is the set of accessible qubit pointers, and where $l\in\mathcal{L}(\mathbb{N})$ is the list of qubit pointers under consideration. Let $\text{Conf}_{n}$, be the set of configurations over $n$ qubits. The initial configuration in $\text{Conf}_{n}$ of a program $\mathrm{D}::\mathrm{S}$ on input state $\left|\psi\right\rangle\in\mathcal{H}_{2^{n}}$ is $c_{init}(\left|\psi\right\rangle)\triangleq(\mathrm{S},\left|\psi\right\rangle%
,\{1,\dots,n\},[1,\dots,n])$. A final configuration can be defined in the same way as $c_{final}(\left|\psi\right\rangle)\triangleq(\top,\left|\psi\right\rangle,\{1,%
\dots,n\},[1,\dots,n])$.\par
 
% END S2.SS2.p4
 
% BEGIN S2.SS2.p5
 
Each unitary transformation $\mathrm{U}$ of a unitary application $\bar{\mathrm{q}}\textnormal{\small[$\mathrm{i}$]}{\ \textnormal{{%
\textasteriskcentered=}}\ }\mathrm{U}^{f}(\mathrm{j});$, comes with a function $\llbracket\mathrm{U}\rrbracket$ assigning a unitary matrix $\llbracket\mathrm{U}\rrbracket(f)(n)\in\mathbb{C}^{2\times 2}$ to each integer $n$ and polynomial-time approximable total function $f\in\mathbb{Z}\to[0,2\pi)$. For example, the gates of the quantum Fourier transform can be defined by $R_{n}\triangleq\llbracket\textnormal{Ph}\rrbracket(\lambda x.\pi/2^{x-1})(n)$ with $\llbracket\textnormal{Ph}\rrbracket(f)(n)\triangleq(\begin{smallmatrix}1&0\\
0&e^{if(n)}\end{smallmatrix})$. The other basic unary gates are the NOT and the $R_{Y}$ gate.\par
 
% END S2.SS2.p5
 
% BEGIN S2.SS2.p6
 
The big-step semantics $\cdot\stackrel{{\scriptstyle\cdot}}{{\longrightarrow}}\cdot$ is defined in Figure 4 as a relation in $\bigcup_{n\in\mathbb{N}}\text{Conf}_{n}\times\mathbb{N}\times\text{Conf}_{n}$. Standard notations from quantum computation are used such as tensor product $\otimes$, $\bra{\psi}$ for the conjugate transpose of $\left|\psi\right\rangle$, or given a dimension $m$, $\left|k\right\rangle_{n}\triangleq I_{2^{n-1}}\otimes\left|k\right\rangle%
\otimes I_{2^{m-n}}$ for $k\in\{0,1\}$. In Figure 4, the set $A$ of accessible qubits is used to ensure that unitary operations on qubits can be physically implemented. For example, to ensure reversibility, in a quantum branch $\text{\color[rgb]{0.03515625,0.0390625,0.72265625}\definecolor[named]{%
pgfstrokecolor}{rgb}{0.03515625,0.0390625,0.72265625}{{qcase }}}{\mathrm{s}%
\textnormal{\small[$\mathrm{i}$]}}\text{\color[rgb]{%
0.03515625,0.0390625,0.72265625}\definecolor[named]{pgfstrokecolor}{rgb}{%
0.03515625,0.0390625,0.72265625}{{ of }}}\{0\to\mathrm{S}_{0},\,1\to\mathrm{S}%
_{1}\}$, statements $\mathrm{S}_{0}$ and $\mathrm{S}_{1}$ cannot access $\mathrm{s}\textnormal{\small[$\mathrm{i}$]}$.\par
 
% END S2.SS2.p6
 
% BEGIN Thmtheorem1
 
\par\medskip\noindent\textbf{Definition 1.}\par
 
% BEGIN Thmtheorem1.p1
 
The time complexity $\textnormal{Time}_{\mathrm{P}}:\mathbb{N}\to\mathbb{N}$ of a program $\mathrm{P}\triangleq\mathrm{D}::\mathrm{S}$ is defined by\par
    \[\textnormal{Time}_{\mathrm{P}}(n)\triangleq\max_{\left|\phi\right\rangle\in%
\mathcal{H}_{2^{n}}}\{m\in\mathbb{N}\mid\exists\left|\phi^{\prime}\right%
\rangle\in\mathcal{H}_{2^{n}},\ c_{init}(\left|\phi\right\rangle)\stackrel{{%
\scriptstyle m}}{{\longrightarrow}}c_{final}(\left|\phi^{\prime}\right\rangle)\}.\]    
% END Thmtheorem1.p1
 
% END Thmtheorem1
 
% BEGIN S2.SS2.p7
 
Intuitively, when $c\stackrel{{\scriptstyle m}}{{\longrightarrow}}c^{\prime}$ holds, the time complexity $m$ is an integer corresponding to the maximum number of procedure calls performed over each (classical and quantum) branch during the evaluation of a configuration $c\in\text{Conf}_{n}$. We write $\llbracket\mathrm{P}\rrbracket(\left|\psi\right\rangle)=\left|\psi^{\prime}\right\rangle$, whenever $c_{init}(\left|\psi\right\rangle)\stackrel{{\scriptstyle m}}{{\longrightarrow}%
}c_{final}(\left|\psi^{\prime}\right\rangle)$ holds for some $m$. If the program $\mathrm{P}$ terminates on any input (i.e., always reaches a final configuration) then $\llbracket\mathrm{P}\rrbracket$ is a total function on quantum states.\par
 
% END S2.SS2.p7
 
% BEGIN Thmtheorem2
 
\par\medskip\noindent\textbf{Example 2.}\par
 
% BEGIN Thmtheorem2.p1
 
For the program PAIRS of Figure 1, $\textnormal{Time}_{\texttt{PAIRS}}(n)=\lfloor\frac{n}{2}\rfloor+1$ since each procedure call removes two qubits until it reaches a sorted set $\bar{\mathrm{q}}$ such that $|\bar{\mathrm{q}}|\leq 1$ (depending on whether $n$ is odd or even) and for both sizes there are no more procedure calls.\par
 
% END Thmtheorem2.p1
 
% END Thmtheorem2
 
% END S2.SS2
 
% BEGIN S2.SS3
 
\par\medskip\noindent\textbf{2.3 Polynomial Branching Programs}\par
 
% BEGIN S2.F6

\begin{center}\includegraphics[width=0.98\linewidth,height=0.70\textheight,keepaspectratio]{graphic-05.pdf}\end{center}
 
Figure 6: Program QFT for the quantum Fourier transform.\par
 
% END S2.F6
 
% BEGIN S2.SS3.p1
 
We define three restrictions on the programming language to consider only a strict subset, called pbp for Polynomial Branching Programs, on which branch sequentialization can be avoided. First, we define a well-foundedness criterion to consider only terminating programs. A program $\mathrm{P}$ is well-founded if each recursive procedure call removes at least one qubit in its parameter. wf denotes the set of well-founded programs. Then, we define a criterion to exclude programs with exponential runtime. Toward that end, the notion of width of a procedure proc in a program $\mathrm{P}$ is introduced.\par
 
% END S2.SS3.p1
 
% BEGIN Thmtheorem3
 
\par\medskip\noindent\textbf{Definition 3.}\par
 
% BEGIN Thmtheorem3.p1
 
Given a program $\mathrm{P}$, the width of a procedure $\text{\color[rgb]{0.79296875,0.3203125,0.8359375}\definecolor[named]{%
pgfstrokecolor}{rgb}{0.79296875,0.3203125,0.8359375}{{proc}}}\in\mathrm{P}$, denoted $0pt_{\mathrm{P}}(\text{\color[rgb]{0.79296875,0.3203125,0.8359375}\definecolor%
[named]{pgfstrokecolor}{rgb}{0.79296875,0.3203125,0.8359375}{{proc}}})$, is defined as $0pt_{\mathrm{P}}(\text{\color[rgb]{0.79296875,0.3203125,0.8359375}\definecolor%
[named]{pgfstrokecolor}{rgb}{0.79296875,0.3203125,0.8359375}{{proc}}})%
\triangleq w^{\text{\color[rgb]{0.79296875,0.3203125,0.8359375}\definecolor[%
named]{pgfstrokecolor}{rgb}{0.79296875,0.3203125,0.8359375}{{proc}}}}_{\mathrm%
{P}}(\mathrm{S}^{\text{\color[rgb]{0.79296875,0.3203125,0.8359375}\definecolor%
[named]{pgfstrokecolor}{rgb}{0.79296875,0.3203125,0.8359375}{{proc}}}})$, where $w_{\mathrm{P}}^{\text{\color[rgb]{0.79296875,0.3203125,0.8359375}\definecolor[%
named]{pgfstrokecolor}{rgb}{0.79296875,0.3203125,0.8359375}{{proc}}}}(\mathrm{%
S})$ is defined inductively as follows:\par
    $\displaystyle w_{\mathrm{P}}^{\text{\color[rgb]{0.79296875,0.3203125,0.8359375%
}\definecolor[named]{pgfstrokecolor}{rgb}{0.79296875,0.3203125,0.8359375}{{%
proc}}}}(\text{\color[rgb]{0.03515625,0.0390625,0.72265625}\definecolor[named]%
{pgfstrokecolor}{rgb}{0.03515625,0.0390625,0.72265625}{{skip}}};)$ $\displaystyle\triangleq 0$     $\displaystyle w_{\mathrm{P}}^{\text{\color[rgb]{0.79296875,0.3203125,0.8359375%
}\definecolor[named]{pgfstrokecolor}{rgb}{0.79296875,0.3203125,0.8359375}{{%
proc}}}}(\mathrm{q}\emph{{\ {{\textasteriskcentered=}}\ }}\mathrm{U}^{f}(%
\mathrm{i});)$ $\displaystyle\triangleq 0,$     $\displaystyle w_{\mathrm{P}}^{\text{\color[rgb]{0.79296875,0.3203125,0.8359375%
}\definecolor[named]{pgfstrokecolor}{rgb}{0.79296875,0.3203125,0.8359375}{{%
proc}}}}(\mathrm{S}_{1}\ \mathrm{S}_{2})$ $\displaystyle\triangleq w^{\text{\color[rgb]{0.79296875,0.3203125,0.8359375}%
\definecolor[named]{pgfstrokecolor}{rgb}{0.79296875,0.3203125,0.8359375}{{proc%
}}}}_{\mathrm{P}}(\mathrm{S}_{1})+w^{\text{\color[rgb]{%
0.79296875,0.3203125,0.8359375}\definecolor[named]{pgfstrokecolor}{rgb}{%
0.79296875,0.3203125,0.8359375}{{proc}}}}_{\mathrm{P}}(\mathrm{S}_{2}),$     $\displaystyle w_{\mathrm{P}}^{\text{\color[rgb]{0.79296875,0.3203125,0.8359375%
}\definecolor[named]{pgfstrokecolor}{rgb}{0.79296875,0.3203125,0.8359375}{{%
proc}}}}(\text{\color[rgb]{0.03515625,0.0390625,0.72265625}\definecolor[named]%
{pgfstrokecolor}{rgb}{0.03515625,0.0390625,0.72265625}{{if }}}\mathrm{b}\text{%
\color[rgb]{0.03515625,0.0390625,0.72265625}\definecolor[named]{pgfstrokecolor%
}{rgb}{0.03515625,0.0390625,0.72265625}{{ then }}}\mathrm{S}_{0}\text{\color[%
rgb]{0.03515625,0.0390625,0.72265625}\definecolor[named]{pgfstrokecolor}{rgb}{%
0.03515625,0.0390625,0.72265625}{{ else }}}\mathrm{S}_{1})$ $\displaystyle\triangleq\max(w^{\text{\color[rgb]{%
0.79296875,0.3203125,0.8359375}\definecolor[named]{pgfstrokecolor}{rgb}{%
0.79296875,0.3203125,0.8359375}{{proc}}}}_{\mathrm{P}}(\mathrm{S}_{0}),w^{%
\text{\color[rgb]{0.79296875,0.3203125,0.8359375}\definecolor[named]{%
pgfstrokecolor}{rgb}{0.79296875,0.3203125,0.8359375}{{proc}}}}_{\mathrm{P}}(%
\mathrm{S}_{1})),$     $\displaystyle w_{\mathrm{P}}^{\text{\color[rgb]{0.79296875,0.3203125,0.8359375%
}\definecolor[named]{pgfstrokecolor}{rgb}{0.79296875,0.3203125,0.8359375}{{%
proc}}}}(\text{\color[rgb]{0.03515625,0.0390625,0.72265625}\definecolor[named]%
{pgfstrokecolor}{rgb}{0.03515625,0.0390625,0.72265625}{{qcase }}}{\mathrm{q}}%
\text{\color[rgb]{0.03515625,0.0390625,0.72265625}\definecolor[named]{%
pgfstrokecolor}{rgb}{0.03515625,0.0390625,0.72265625}{{ of }}}\{0\to\mathrm{S}%
_{0},\,1\to\mathrm{S}_{1}\})$ $\displaystyle\triangleq\max(w^{\text{\color[rgb]{%
0.79296875,0.3203125,0.8359375}\definecolor[named]{pgfstrokecolor}{rgb}{%
0.79296875,0.3203125,0.8359375}{{proc}}}}_{\mathrm{P}}(\mathrm{S}_{0}),w^{%
\text{\color[rgb]{0.79296875,0.3203125,0.8359375}\definecolor[named]{%
pgfstrokecolor}{rgb}{0.79296875,0.3203125,0.8359375}{{proc}}}}_{\mathrm{P}}(%
\mathrm{S}_{1})),$     $\displaystyle w_{\mathrm{P}}^{\text{\color[rgb]{0.79296875,0.3203125,0.8359375%
}\definecolor[named]{pgfstrokecolor}{rgb}{0.79296875,0.3203125,0.8359375}{{%
proc}}}}(\text{\color[rgb]{0.03515625,0.0390625,0.72265625}\definecolor[named]%
{pgfstrokecolor}{rgb}{0.03515625,0.0390625,0.72265625}{{call }}}\text{\color[%
rgb]{0.79296875,0.3203125,0.8359375}\definecolor[named]{pgfstrokecolor}{rgb}{%
0.79296875,0.3203125,0.8359375}{{proc'}}}(\mathrm{s});)$ $\displaystyle\triangleq\begin{cases}1&\text{if }\text{\color[rgb]{%
0.79296875,0.3203125,0.8359375}\definecolor[named]{pgfstrokecolor}{rgb}{%
0.79296875,0.3203125,0.8359375}{{proc}}}\sim_{\mathrm{P}}\text{\color[rgb]{%
0.79296875,0.3203125,0.8359375}\definecolor[named]{pgfstrokecolor}{rgb}{%
0.79296875,0.3203125,0.8359375}{{proc'}}},\\
0&\text{otherwise}.\end{cases}$    
Let $\textnormal{{width}}_{\leq 1}$ be the set of programs with procedures of width at most $1$.\par
 
% END Thmtheorem3.p1
 
% END Thmtheorem3
 
% BEGIN S2.SS3.p2
 
Programs of width $1$ are inherently polynomial as they cannot perform an exponential number of procedure calls in sequence. However, the total number of calls in superposition (that is to say, across all branches of computation) may be exponential for such programs by the definition of the width for conditional and quantum case. For this reason, a compilation strategy whose circuits scale as the sum of the branches instead of their maximum can lead to circuits of exponential size.\par
 
% END S2.SS3.p2
 
% BEGIN S2.SS3.p3
 
Finally, for the purpose of our compilation process, we impose further syntactical conditions on programs. We restrict our attention to what we will call basic programs.\par
 
% END S2.SS3.p3
 
% BEGIN Thmtheorem4
 
\par\medskip\noindent\textbf{Definition 4.}\par
 
% BEGIN Thmtheorem4.p1
 
A program is basic if there exists a sorted set $\mathrm{s}$ such that every procedure call is of the shape $\text{\color[rgb]{0.03515625,0.0390625,0.72265625}\definecolor[named]{%
pgfstrokecolor}{rgb}{0.03515625,0.0390625,0.72265625}{{call}}}\ \text{\color[%
rgb]{0.79296875,0.3203125,0.8359375}\definecolor[named]{pgfstrokecolor}{rgb}{%
0.79296875,0.3203125,0.8359375}{{proc}}}(\bar{\mathrm{q}})$ or $\text{\color[rgb]{0.03515625,0.0390625,0.72265625}\definecolor[named]{%
pgfstrokecolor}{rgb}{0.03515625,0.0390625,0.72265625}{{call}}}\ \text{\color[%
rgb]{0.79296875,0.3203125,0.8359375}\definecolor[named]{pgfstrokecolor}{rgb}{%
0.79296875,0.3203125,0.8359375}{{proc}}}(\mathrm{s})$. Let basic denote the set of basic programs.\par
 
% END Thmtheorem4.p1
 
% END Thmtheorem4
 
% BEGIN S2.SS3.p4
 
basic programs are programs whose procedure calls are restricted to the full qubit list or to a fixed qubit list expression. This restriction simplifies the treatment of recursive procedure calls during compilation.\par
 
% END S2.SS3.p4
 
% BEGIN Thmtheorem5
 
\par\medskip\noindent\textbf{Definition 5 (Polynomially Branching Programs).}\par
 
% BEGIN Thmtheorem5.p1
 
The set pbp of polynomially branching programs is defined as $\textnormal{{pbp}}{}\triangleq\textnormal{{wf}}\cap\textnormal{{width}}_{\leq 1%
}\cap\textnormal{{basic}}{}$.\par
 
% END Thmtheorem5.p1
 
% END Thmtheorem5
 
% BEGIN S2.SS3.p5
 
Notice that pbp is strictly included in the programming language of [17], that roughly corresponds to $\textnormal{{wf}}\cap\textnormal{{width}}_{\leq 1}$, where procedure calls can take an extra integer parameter. Since the language of [17] is polytime-sound, so is pbp (Theorem 11). The restriction to programs in pbp does not compromise the extensional completeness of the language as, by Theorem 11, pbp is complete for the class fbqp of functions in $\{0,1\}^{*}\to\{0,1\}^{*}$, i.e., functions that can be approximated with at least $2/3$ probability by a quantum Turing machine running in polynomial time [6, 29].\par
 
% END S2.SS3.p5
 
% BEGIN Thmtheorem6
 
\par\medskip\noindent\textbf{Example 6.}\par
 
% BEGIN Thmtheorem6.p1
 
The QFT program written in Figure 6 is in pbp. Indeed, the program QFT is in basic as calls are only performed on either $\bar{\mathrm{q}}$ or $\bar{\mathrm{q}}\ominus\textnormal{\small[$-1$]}$. It is also in wf as all recursive calls are performed on parameter $\bar{\mathrm{q}}\ominus\textnormal{\small[$-1$]}$. Finally, it is in $\textnormal{{width}}_{\leq 1}$ as $0pt_{\texttt{QFT}}(\text{\color[rgb]{0.79296875,0.3203125,0.8359375}%
\definecolor[named]{pgfstrokecolor}{rgb}{0.79296875,0.3203125,0.8359375}{{qft}%
}})=0pt_{\texttt{QFT}}(\text{\color[rgb]{0.79296875,0.3203125,0.8359375}%
\definecolor[named]{pgfstrokecolor}{rgb}{0.79296875,0.3203125,0.8359375}{{rot}%
}})=0pt_{\texttt{QFT}}(\text{\color[rgb]{0.79296875,0.3203125,0.8359375}%
\definecolor[named]{pgfstrokecolor}{rgb}{0.79296875,0.3203125,0.8359375}{{%
shift}}})=1$.\par
 
% END Thmtheorem6.p1
 
% END Thmtheorem6
 
% END S2.SS3
 
% END S2


% BEGIN S3
 
\par\medskip\noindent\textbf{3 Compilation Strategy}\par
 
% BEGIN S3.p1
 
We now present the compilation algorithm from $\textnormal{{wf}}\cap\textnormal{{width}}_{\leq 1}$ to quantum circuits based on the merging strategy of Figure 2(c). Compilation is restricted to programs in $\textnormal{{wf}}\cap\textnormal{{width}}_{\leq 1}$ for two reasons. The well-foundedness criterion wf ensures that the compilation always terminates. The restriction to $\textnormal{{width}}_{\leq 1}$ prevents exponential blow-up. Note however that, in Theorem 9, the extra basic restriction is required to avoid the branch sequentialization problem, that is, branch sequentialization is avoided on pbp.\par
 
% END S3.p1
 
% BEGIN S3.SS1
 
\par\medskip\noindent\textbf{3.1 Anchoring and Merging}\par
 
% BEGIN S3.F7

\begin{center}\includegraphics[width=0.98\linewidth,height=0.70\textheight,keepaspectratio]{graphic-06.pdf}\end{center}
 
Figure 7: Anchoring and merging.\par
 
% END S3.F7
 
% BEGIN S3.SS1.p1
 
A control structure $cs$ is a partial map in $\mathbb{N}\to\{0,1\}$. Intuitively, $cs$ represents a map from qubit pointers to their control values for a controlled gate and will also be used as a shorthand notation in circuits. For $n\in\mathbb{N}$ and $k\in\{0,1\}$, let $cs[n:=k]$ be the control structure obtained from $cs$ by setting $cs(n)=k$. We denote by $dom(cs)$ the domain of the control structure and by $\cdot$, the control structure such that $dom(\cdot)=\emptyset$.\par
 
% END S3.SS1.p1
 
% BEGIN S3.SS1.p2
 
During the compilation of the program statement, every recursive call is handled as follows: if it is the first call with this procedure name and input size, an ancilla is created and its procedure statement is compiled under a control on this ancilla (anchoring), otherwise this procedure call has already been compiled and the control structure of the new call is sent into the corresponding ancilla (merging). Anchoring and merging are represented in Figure 7 as rewrite rules on quantum circuits for a procedure call $\text{\color[rgb]{0.03515625,0.0390625,0.72265625}\definecolor[named]{%
pgfstrokecolor}{rgb}{0.03515625,0.0390625,0.72265625}{{call }}}\text{\color[%
rgb]{0.79296875,0.3203125,0.8359375}\definecolor[named]{pgfstrokecolor}{rgb}{%
0.79296875,0.3203125,0.8359375}{{proc}}}(\mathrm{s})$, $cs$ denoting the control structure corresponding to the (possibly nested) quantum cases, where the procedure call occurs.\par
 
% END S3.SS1.p2
 
% BEGIN S3.SS1.p3
 
$a_{k}^{\text{\color[rgb]{0.79296875,0.3203125,0.8359375}\definecolor[named]{%
pgfstrokecolor}{rgb}{0.79296875,0.3203125,0.8359375}{{proc}}}}$ represents the ancilla introduced (in case of anchoring) or reused (in case of merging). The integer $k$ refers to the number $|\mathrm{s}|$ of qubits in the procedure call. The grey box materializes the controlled statements left to compile.\par
 
% END S3.SS1.p3
 
% END S3.SS1
 
% BEGIN S3.SS2
 
\par\medskip\noindent\textbf{3.2 The compile Algorithm}\par
 
% BEGIN S3.F8

\begin{center}\includegraphics[width=0.98\linewidth,height=0.70\textheight,keepaspectratio]{graphic-07.pdf}\end{center}
 
Figure 8: Rewrite rules of compile.\par
 
% END S3.F8
 
% BEGIN S3.SS2.p1
 
Let controlled statements be pairs of the shape $(cs,\mathrm{S})$, for some control structure $cs$ and some statement $\mathrm{S}$. In the compilation process, controlled statements $(cs,\mathrm{S})$ are used to represent a statement $\mathrm{S}$ that remains to be compiled into a quantum circuit $C$ together with a control structure $cs$, representing the qubits controlling the compiled circuit $C$. The compilation algorithm compile is described by the rewrite rules of Figure 8. Generating the circuit corresponding to the program $\mathrm{P}=\mathrm{D}::\mathrm{S}$ will consist in running compile on the controlled statement $(\cdot,\mathrm{S})$ for a fixed list of qubits pointers $[1,\ldots,n]$ (hence, an input of fixed size $n$). We denote by $\textnormal{{compile}}{}(\mathrm{P},n)$ the circuit obtained for program $\mathrm{P}$ on size $n$. The algorithm standardly generates the quantum circuit corresponding to a $\textnormal{{wf}}\cap\textnormal{{width}}_{\leq 1}$ program inductively on the statement $\mathrm{S}$. Indeed, in the rules of Figure 8 when compile is called on a given controlled statement $(cs,\mathrm{S})$, an inductive call to compile is performed on controlled statements whose statements are sub-statements of $\mathrm{S}$. The two base cases are the rules for the skip statement and the unitary application. In these cases, the compilation just outputs the identity circuit and a controlled gate computing the unitary, respectively. The rules for sequence and quantum case perform two inductive calls to compile on each branch of the statement. The rule for quantum case is the only rule that directly performs changes on the control structure. In the particular case of a call to a recursive procedure (i.e., when $w_{\mathrm{P}}^{\text{\color[rgb]{0.79296875,0.3203125,0.8359375}\definecolor[%
named]{pgfstrokecolor}{rgb}{0.79296875,0.3203125,0.8359375}{{proc}}}}(\mathrm{%
S}^{\text{\color[rgb]{0.79296875,0.3203125,0.8359375}\definecolor[named]{%
pgfstrokecolor}{rgb}{0.79296875,0.3203125,0.8359375}{{proc}}}})>0$), compile calls the optimize subroutine to perform anchoring and merging. This call to optimize is highlighted in Figure 8 through the use of a shaded square \raisebox{-0.15em}{\includegraphics[height=1.1em]{x78.png}}, which takes the procedure name proc as superscript. We call this process the optimization of procedure proc.\par
 
% END S3.SS2.p1
 
% BEGIN S3.F9
   
% BEGIN S3.F9.fig1

\begin{center}\includegraphics[width=0.98\linewidth,height=0.70\textheight,keepaspectratio]{graphic-08.pdf}\end{center}
 
% END S3.F9.fig1
    
% BEGIN S3.F9.fig2
 
\begin{center}\includegraphics[width=0.98\linewidth,height=0.70\textheight,keepaspectratio]{graphic-09.pdf}\end{center}
 
% END S3.F9.fig2
    
% BEGIN S3.F9.fig3
 
\begin{center}\includegraphics[width=0.98\linewidth,height=0.70\textheight,keepaspectratio]{graphic-10.pdf}\end{center}
 
% END S3.F9.fig3
    
% BEGIN S3.F9.fig4
 
\begin{center}\includegraphics[width=0.98\linewidth,height=0.70\textheight,keepaspectratio]{graphic-11.pdf}\end{center}
 
% END S3.F9.fig4
    
% BEGIN S3.F9.fig5
 
\begin{center}\includegraphics[width=0.98\linewidth,height=0.70\textheight,keepaspectratio]{graphic-12.pdf}\end{center}
 
% END S3.F9.fig5
    
% BEGIN S3.F9.fig6
 
\begin{center}\includegraphics[width=0.98\linewidth,height=0.70\textheight,keepaspectratio]{graphic-13.pdf}\end{center}
 
% END S3.F9.fig6
   
Figure 9: Rewrite rules of optimize.\par
 
% END S3.F9
 
% BEGIN S3.SS2.p2
 
The rewrite rules for the subroutine optimize on procedure proc are described in Figure 9. This subroutine takes two input parameters:\par
  
\par\noindent  $\blacksquare$ 
% BEGIN S3.I1.i1.p1
 
a first list $l_{\mathrm{Cst}}$ of controlled statements, the ones to be optimized,\par
 
% END S3.I1.i1.p1
 
 
\par\noindent  $\blacksquare$ 
% BEGIN S3.I1.i2.p1
 
a second list of controlled statements, called contextual list called $l_{\text{M}}$ and denoted by \raisebox{-0.15em}{\includegraphics[height=1.1em]{x79.png}}, consisting in controlled statements that are orthogonal to those in $l_{\mathrm{Cst}}$ and do not contain recursive procedure calls.\par
 
% END S3.I1.i2.p1
 
  
As described by the last rule of Figure 8, each initial call to optimize is performed on the singleton list containing one controlled statement $(cs,\mathrm{S}^{\text{\color[rgb]{0.79296875,0.3203125,0.8359375}\definecolor[%
named]{pgfstrokecolor}{rgb}{0.79296875,0.3203125,0.8359375}{{proc}}}}\{\mathrm%
{s}/\bar{\mathrm{q}}\})$ and a contextual circuit equal to the identity circuit.\par
 
% END S3.SS2.p2
 
% BEGIN S3.SS2.p3
 
In Figure 9, the rules are applied by considering the first controlled statement in the list $l_{\mathrm{Cst}}$. We just specify the most interesting case below:\par
  
\par\noindent  $\blacksquare$ 
% BEGIN S3.I2.i1.p1
 
For a controlled statement $(cs,\mathrm{S}_{1}\ \mathrm{S}_{2})$, two distinct rules can be applied. In the first scenario, where $w_{\mathrm{P}}^{\text{\color[rgb]{0.79296875,0.3203125,0.8359375}\definecolor[%
named]{pgfstrokecolor}{rgb}{0.79296875,0.3203125,0.8359375}{{proc}}}}(\mathrm{%
S}_{1})=1$, we have that $\mathrm{S}_{1}$ contains a recursive procedure call and $\mathrm{S}_{2}$ does not (this is a consequence of the $\textnormal{{width}}_{\leq 1}$ condition in pbp), or the converse. Depending on this, the rule select on which control statement $(cs,\mathrm{S}_{1})$ or $(cs,\mathrm{S}_{2})$ to perform the optimization and, append the compiled circuit of the other controlled statement to the left or to the right.\par
 
% END S3.I2.i1.p1
 
 
\par\noindent  $\blacksquare$ 
% BEGIN S3.I2.i2.p1
 
For a controlled statement with an if statement, we precompute the boolean value $b$ (which only depends on classical values such as the integer input and the register size), and according to whether the corresponding branch contains a recursive call or not, we either add the controlled statement to $l_{\mathrm{Cst}}$ or the contextual list.\par
 
% END S3.I2.i2.p1
 
 
\par\noindent  $\blacksquare$ 
% BEGIN S3.I2.i3.p1
 
For a controlled statement with a qcase statement, one or two of the branches are added to the list $l_{\mathrm{Cst}}$ depending on whether they are both recursive or not. The branch that is not recursive is added to the contextual list.\par
 
% END S3.I2.i3.p1
 
 
\par\noindent  $\blacksquare$ 
% BEGIN S3.I2.i4.p1
 
For a controlled statement with a procedure call $\text{\color[rgb]{0.03515625,0.0390625,0.72265625}\definecolor[named]{%
pgfstrokecolor}{rgb}{0.03515625,0.0390625,0.72265625}{{call }}}\text{\color[%
rgb]{0.79296875,0.3203125,0.8359375}\definecolor[named]{pgfstrokecolor}{rgb}{%
0.79296875,0.3203125,0.8359375}{{proc'}}}(\mathrm{s});$ we perform either anchoring or merging depending on whether the ancilla $a_{|\mathrm{s}|}^{\text{\color[rgb]{0.79296875,0.3203125,0.8359375}%
\definecolor[named]{pgfstrokecolor}{rgb}{0.79296875,0.3203125,0.8359375}{{proc%
'}}}}$ exists or not. In the case of merging, we include a controlled permutation of qubit addresses (as described in Lemma 12) in order to match the input qubits of the different procedures. Notice that, for basic programs, the input qubits will match and therefore this permutation is trivial, and merging is as described in Figure 7.\par
 
% END S3.I2.i4.p1
 
  
% END S3.SS2.p3
 
% BEGIN S3.SS2.p4
 
At the end of optimize, when $l_{\mathrm{Cst}}=[]$, we rearrange the contextual list in the following way. Let $[(cs_{1},\mathrm{S}_{1},l_{1}),\dots,(cs_{k},\mathrm{S}_{k},l_{k})]$ be the state of the contextual list at the end of optimize. We may rewrite each controlled statement as a list of its atomic elements, with a transformation denoted seq. Let $(\mathrm{q},l)\Downarrow_{\mathbb{N}}n$, then\par
    $\displaystyle\text{seq}(cs,\text{\color[rgb]{0.03515625,0.0390625,0.72265625}%
\definecolor[named]{pgfstrokecolor}{rgb}{0.03515625,0.0390625,0.72265625}{{%
skip}}};,l)$ $\displaystyle\triangleq\ [],$     $\displaystyle\text{seq}(cs,\mathrm{q}{\ \textnormal{{\textasteriskcentered=}}%
\ }\text{U}^{f}(\mathrm{j});,l)$ $\displaystyle\triangleq\ [(cs,\mathrm{q}{\ \textnormal{{\textasteriskcentered=%
}}\ }\text{U}^{f}(\mathrm{j});,l)],$     $\displaystyle\text{seq}(cs,\mathrm{S}_{1}\ \mathrm{S}_{2},l)$ $\displaystyle\triangleq\ \text{seq}(cs,\mathrm{S}_{1},l)@\text{seq}(cs,\mathrm%
{S}_{1},l),$     $\displaystyle\text{seq}(cs,\text{\color[rgb]{0.03515625,0.0390625,0.72265625}%
\definecolor[named]{pgfstrokecolor}{rgb}{0.03515625,0.0390625,0.72265625}{{if %
}}}\mathrm{b}\text{\color[rgb]{0.03515625,0.0390625,0.72265625}\definecolor[%
named]{pgfstrokecolor}{rgb}{0.03515625,0.0390625,0.72265625}{{ then }}}\mathrm%
{S}_{\text{\color[rgb]{0.03515625,0.0390625,0.72265625}\definecolor[named]{%
pgfstrokecolor}{rgb}{0.03515625,0.0390625,0.72265625}{{true}}}}\text{\color[%
rgb]{0.03515625,0.0390625,0.72265625}\definecolor[named]{pgfstrokecolor}{rgb}{%
0.03515625,0.0390625,0.72265625}{{ else }}}\mathrm{S}_{\text{\color[rgb]{%
0.03515625,0.0390625,0.72265625}\definecolor[named]{pgfstrokecolor}{rgb}{%
0.03515625,0.0390625,0.72265625}{{false}}}},l)$ $\displaystyle\triangleq\ \text{seq}(cs,\mathrm{S}_{b},l),\text{ if }(\mathrm{b%
},l)\Downarrow_{\mathbb{B}}b,$     $\displaystyle\text{seq}(cs,\text{\color[rgb]{0.03515625,0.0390625,0.72265625}%
\definecolor[named]{pgfstrokecolor}{rgb}{0.03515625,0.0390625,0.72265625}{{%
qcase }}}\mathrm{q}\text{\color[rgb]{0.03515625,0.0390625,0.72265625}%
\definecolor[named]{pgfstrokecolor}{rgb}{0.03515625,0.0390625,0.72265625}{{ of%
 }}}\{0\to\mathrm{S}_{0},\,1\to\mathrm{S}_{1}\},l)$ $\displaystyle\triangleq\ \text{seq}(cs[n:=0],\mathrm{S}_{0},l)@\text{seq}(cs[n%
:=1],\mathrm{S}_{1},l),$     $\displaystyle\text{seq}(cs,\text{\color[rgb]{0.03515625,0.0390625,0.72265625}%
\definecolor[named]{pgfstrokecolor}{rgb}{0.03515625,0.0390625,0.72265625}{{%
call }}}\text{\color[rgb]{0.79296875,0.3203125,0.8359375}\definecolor[named]{%
pgfstrokecolor}{rgb}{0.79296875,0.3203125,0.8359375}{{proc}}}(\mathrm{s});,l)$ $\displaystyle\triangleq[(cs,\text{\color[rgb]{0.03515625,0.0390625,0.72265625}%
\definecolor[named]{pgfstrokecolor}{rgb}{0.03515625,0.0390625,0.72265625}{{%
call }}}\text{\color[rgb]{0.79296875,0.3203125,0.8359375}\definecolor[named]{%
pgfstrokecolor}{rgb}{0.79296875,0.3203125,0.8359375}{{proc}}}(\mathrm{s});,l)].$    
% END S3.SS2.p4
 
% BEGIN S3.SS2.p5
 
This separation of statements allows for a partitioning according to the type of procedure call appearing in the statement. Given a list of controlled statements $\mathcal{L}$, we define the list $[\mathcal{L}_{0},\mathcal{L}_{1},\dots,\mathcal{L}_{m}]$ where, for procedures proc1, …, procm that are not mutually recursive.\par
    $\displaystyle\mathcal{L}_{0}$ $\displaystyle\triangleq\{(cs,\mathrm{S},l)\in\mathcal{L}:\not\exists\text{%
\color[rgb]{0.79296875,0.3203125,0.8359375}\definecolor[named]{pgfstrokecolor}%
{rgb}{0.79296875,0.3203125,0.8359375}{{proc}}}\text{ such that }w_{\mathrm{P}}%
^{\text{\color[rgb]{0.79296875,0.3203125,0.8359375}\definecolor[named]{%
pgfstrokecolor}{rgb}{0.79296875,0.3203125,0.8359375}{{proc}}}}(\mathrm{S}^{%
\text{\color[rgb]{0.79296875,0.3203125,0.8359375}\definecolor[named]{%
pgfstrokecolor}{rgb}{0.79296875,0.3203125,0.8359375}{{proc}}}})=1\text{ and }w%
_{\mathrm{P}}^{\text{\color[rgb]{0.79296875,0.3203125,0.8359375}\definecolor[%
named]{pgfstrokecolor}{rgb}{0.79296875,0.3203125,0.8359375}{{proc}}}}(\mathrm{%
S})=1\}.$     $\displaystyle\mathcal{L}_{\text{\color[rgb]{0.79296875,0.3203125,0.8359375}%
\definecolor[named]{pgfstrokecolor}{rgb}{0.79296875,0.3203125,0.8359375}{{proc%
${}_{i}$}}}}$ $\displaystyle\triangleq\{(cs,\mathrm{S},l)\in\mathcal{L}:w_{\mathrm{P}}^{\text%
{\color[rgb]{0.79296875,0.3203125,0.8359375}\definecolor[named]{pgfstrokecolor%
}{rgb}{0.79296875,0.3203125,0.8359375}{{proc${}_{i}$}}}}(\mathrm{S})=1\},\quad
i%
=1,\dots,m.$    
% END S3.SS2.p5
 
% BEGIN S3.SS2.p6
 
Given our choice of procedures and the controlled statements obtained from seq, this partition is well-defined. Performing these two partitions (first in terms of sequential order of statements, and then according to procedure calls), we are able to rewrite:\par
 
% END S3.SS2.p6
 
% BEGIN S3.SS2.p7
 
\begin{center}\includegraphics[width=0.98\linewidth,height=0.70\textheight,keepaspectratio]{graphic-14.pdf}\end{center}
 
% END S3.SS2.p7
 
% BEGIN S3.SS2.p8
 
The different instances of optimize contain calls to procedures that are mutually recursive, which will allow for further anchoring and merging.\par
 
% END S3.SS2.p8
 
% END S3.SS2
 
% BEGIN S3.SS3
 
\par\medskip\noindent\textbf{3.3 Soundness of the algorithm}\par
 
% BEGIN S3.SS3.p1
 
In this section, we discuss the validity of the compilation algorithm. One first observation should be that, given a pbp program, the compilation necessarily terminates. For instance, in compile (Figure 8), all rules besides the procedure call rewrite the controlled statement into either a circuit or into instances of compile of smaller statements. In the case of a procedure call, the rewriting of the procedure body produces a finite number of calls to procedures of lower recursive level.\par
 
% END S3.SS3.p1
 
% BEGIN S3.SS3.p2
 
In optimize, a recursive procedure will result in a finite number of calls to mutually recursive procedures – this is ensured by the well-foundedness condition wf, that requires that recursive procedure calls reduce the size of the input, therefore procedure calls either reduce the level of recursion or the number of available qubits.\par
 
% END S3.SS3.p2
 
% BEGIN S3.SS3.p3
 
The soundness of the compilation algorithm is ensured by an orthogonality invariant in optimize. Let $cs,\,cs^{\prime}$ be two control structures. We say that $cs$ and $cs^{\prime}$ are orthogonal if there exists $i\in\text{dom}(cs)\cap\text{dom}(cs^{\prime})$ such that $cs(i)=1-cs^{\prime}(i)$. Two controlled statements are orthogonal if their control structures are orthogonal.\par
 
% END S3.SS3.p3
 
% BEGIN Thmtheorem7
 
\par\medskip\noindent\textbf{Lemma 7.}\par
 
% BEGIN Thmtheorem7.p1
 
During the compilation of a program $\mathrm{P}\in\textnormal{{wf}}\cap\textnormal{{width}}_{\leq 1}$, for each optimization of a (recursive) procedure proc, all controlled statements in the union of list $l$ and the contextual list $l_{\text{M}}$ are pairwise orthogonal.\par
 
% END Thmtheorem7.p1
 
% END Thmtheorem7
 
% BEGIN S3.SS3.SSS0.P0.SPx1
 
\par\medskip\noindent\textbf{Proof.}\par
 
% BEGIN S3.SS3.SSS0.P0.SPx1.p1
 
This can be checked by inspection of each rewrite rule in Figure 9, as either the same control structure $cs$ is kept in the list, or it is removed, or we consider two control structures, $cs[n=0]$ and $cs[n=1]$, that preserve the orthogonality relations of $cs$ and are also mutually orthogonal due to the control over qubit $n$. $\hfill\blacktriangleleft$\par
 
% END S3.SS3.SSS0.P0.SPx1.p1
 
% BEGIN S3.SS3.SSS0.P0.SPx1.p2
 
This invariant ensures the validity of optimize, and the soundness of the compilation algorithm. It is also a consequence of the $\textnormal{{width}}_{\leq 1}$ restriction in pbp, as by the definition of width, at most one recursive call may appear per branch of a recursive procedure. This ensures that two recursive calls on a given procedure always occur in orthogonal branches and can be simply combined in the same ancilla. Given a circuit $C$, we define its semantics $\llbracket C\rrbracket$ naturally as the composition of the semantics of each gate.\par
 
% END S3.SS3.SSS0.P0.SPx1.p2
 
% BEGIN Thmtheorem8
 
\par\medskip\noindent\textbf{Theorem 8 (Soundness of compilation).}\par
 
% BEGIN Thmtheorem8.p1
 
Given a program $\mathrm{P}\in\textnormal{{wf}}\cap\textnormal{{width}}_{\leq 1}$ and a quantum state $\left|\psi\right\rangle\in\mathcal{H}_{2^{n}}$ we have that $\llbracket\textnormal{{compile}}{}(\mathrm{P},n)\rrbracket(\left|\psi\right%
\rangle)=\llbracket\mathrm{P}\rrbracket(\left|\psi\right\rangle)$.\par
 
% END Thmtheorem8.p1
 
% END Thmtheorem8
 
% END S3.SS3.SSS0.P0.SPx1
 
% END S3.SS3
 
% BEGIN S3.SS4
 
\par\medskip\noindent\textbf{3.4 Illustrating Example}\par
 
% BEGIN S3.F10

\begin{center}\includegraphics[width=0.98\linewidth,height=0.70\textheight,keepaspectratio]{graphic-15.pdf}\end{center}
 
Figure 10: Program REC.\par
 
% END S3.F10
 
% BEGIN S3.SS4.p1
 
We illustrate the compilation process with the program REC, of Figure 10. Notice that REC is in $\textnormal{{wf}}\cap\textnormal{{width}}_{\leq 1}$ but does not belong to pbp, as the procedure rec performs two recursive calls on different sorted sets $\bar{\mathrm{q}}\ominus\textnormal{\small[$1$]}$ and $\bar{\mathrm{q}}\ominus\textnormal{\small[$1,2$]}$. This example thus illustrates that the compilation algorithm is not restricted to pbp.\par
 
% END S3.SS4.p1
 
% BEGIN S3.SS4.p2
 
Given that rec is a recursive procedure, its compilation is performed within optimize. We denote by the empty box \raisebox{-0.15em}{\includegraphics[height=1.1em]{x80.png}} the procedure statement $\mathrm{S}^{\text{\color[rgb]{0.79296875,0.3203125,0.8359375}\definecolor[%
named]{pgfstrokecolor}{rgb}{0.79296875,0.3203125,0.8359375}{{rec}}}}$ and by the dotted box \raisebox{-0.15em}{\includegraphics[height=1.1em]{x81.png}} the statement $\text{\color[rgb]{0.03515625,0.0390625,0.72265625}\definecolor[named]{%
pgfstrokecolor}{rgb}{0.03515625,0.0390625,0.72265625}{{call }}}\text{\color[%
rgb]{0.79296875,0.3203125,0.8359375}\definecolor[named]{pgfstrokecolor}{rgb}{%
0.79296875,0.3203125,0.8359375}{{rec}}}(\mathrm{s});$. Applying rewrite rules to rec we obtain the transitions:\par
 
% END S3.SS4.p2
 
% BEGIN S3.SS4.p3
 
\begin{center}\includegraphics[width=0.98\linewidth,height=0.70\textheight,keepaspectratio]{graphic-16.pdf}\end{center}
 
% END S3.SS4.p3
 
% BEGIN S3.SS4.p4
 
The first step is obtained by applying the rule of if statements, where we assume that $|\bar{\mathrm{q}}|>2$. Afterwards, we apply twice the rule for the qcase statement, adding the two qubits $\bar{\mathrm{q}}\textnormal{\small[$1$]}$ and $\bar{\mathrm{q}}\textnormal{\small[$2$]}$ to the control structure. Finally, the compilation of the skip statement corresponds to the empty circuit. We denote by $\xrightarrow[]{\texttt{rec}}$ the application of this sequence of rewrite rules, which is used in Figure 11.\par
 
% END S3.SS4.p4
 
% BEGIN S3.F11

\begin{center}\includegraphics[width=0.98\linewidth,height=0.70\textheight,keepaspectratio]{graphic-17.pdf}\end{center}
 
Figure 11: Example of anchoring and merging for the program REC in Figure 10.\par
 
% END S3.F11
 
% BEGIN S3.SS4.p5
 
One important thing to notice in Figure 11 is that the compilation strategy does not avoid branch sequentialization locally but rather asymptotically in the construction of the circuit. In other words, the qcase statement in rule rec does generate two instances of $\mathrm{S}^{\text{\color[rgb]{0.79296875,0.3203125,0.8359375}\definecolor[%
named]{pgfstrokecolor}{rgb}{0.79296875,0.3203125,0.8359375}{{rec}}}}$ in the circuit in sequence, one for each branch. However, the merging of calls to rec on inputs of the same size ensures that only one instance per input size needs to be compiled, and therefore this strategy achieves linear complexity in the number of gates.\par
 
% END S3.SS4.p5
 
% END S3.SS4
 
% END S3


% BEGIN S4
 
\par\medskip\noindent\textbf{4 Main Results}\par
 
% BEGIN S4.SS1
 
\par\medskip\noindent\textbf{4.1 No Branch Sequentialization}\par
 
% BEGIN S4.SS1.p1
 
First, we show that the problem of branch sequentialization is solved in pbp. For that purpose, we show that the size of quantum circuit obtained through compiling a pbp program is bounded asymptotically by the time complexity of the program. As the time complexity is the maximum of the complexity of the two branches of a quantum case statement, the branch sequentialization problem is avoided on pbp programs, as the size of the circuit is asymptotically equal to its time semantics. Given a circuit $C$, let $\#C$ denote its size, i.e., number of gates and wires.\par
 
% END S4.SS1.p1
 
% BEGIN Thmtheorem9
 
\par\medskip\noindent\textbf{Theorem 9 (No branch sequentialization).}\par
 
% BEGIN Thmtheorem9.p1
 
Given a program $\mathrm{P}\in\textnormal{{pbp}}{}$ and input $n\in\mathbb{N}$, we have that $\#\textnormal{{{compile}}{}}(\mathrm{P},n)=O(\textnormal{Time}_{\mathrm{P}}(n))$ holds.\par
 
% END Thmtheorem9.p1
 
% END Thmtheorem9
 
% BEGIN S4.SS1.p2
 
One may notice that, in the rules of compile (Figure 8), the compilation of a qcase statement generates two controlled statements in sequence. Similarly, the rule of optimize (Figure 9) for qcase statements appends two controlled statements to the list $l_{\mathrm{Cst}}$ in the case where they are both recursive. This does not pose a problem as, anchoring and merging will ensure that the asymptotic complexity of this statement will be given by the maximum complexity of the branches.\par
 
% END S4.SS1.p2
 
% BEGIN Thmtheorem10
 
\par\medskip\noindent\textbf{Example 10.}\par
 
% BEGIN Thmtheorem10.p1
 
By Theorem 9 and Example 2, we have that $\#\textnormal{{compile}}{}(\texttt{PAIRS},n)=O(n)$.\par
 
% END Thmtheorem10.p1
 
% END Thmtheorem10
 
% END S4.SS1
 
% BEGIN S4.SS2
 
\par\medskip\noindent\textbf{4.2 Polynomial-Time Soundness and Completeness}\par
 
% BEGIN S4.SS2.p1
 
Having shown in Theorem 9 that pbp programs can avoid branch sequentialization, we now turn to the question of whether they constitute an interesting fragment of quantum programs, given that it is restricted by the wf, $\textnormal{{width}}_{\leq 1}$ and basic conditions. We show that the set of pbp programs is sound and complete for quantum polynomial time via an implicit characterization of fbqp, i.e. the set of classical functions that can be approximated with bounded error by a polytime quantum Turing machine [30, 17].\par
 
% END S4.SS2.p1
 
% BEGIN S4.SS2.p2
 
Since the considered programming language does not allow for qubit creation, in order to define functions where the output is larger than the input, we make use of polytime padding. A polytime padding function is a function $f:\{0,1\}^{\ast}\to\{0,1\}^{\ast}$ computable by a (classical) polytime Turing machine such that $f(x)=xy$, for some $y$ depending only on the length of $x$ (and not the value of $x$). Given a set $\mathcal{F}$ of programs, $\llbracket\mathcal{F}\rrbracket$ denotes the set of functions given by $\llbracket\mathrm{P}\rrbracket\circ f$, where $\mathrm{P}\in\mathcal{F}$, $f$ is a polytime padding function, and $\circ$ is the standard function composition. For any $p\in(\frac{1}{2},1]$, we denote by $\llbracket\mathcal{F}\rrbracket_{\geq p}$ the set of functions in $\llbracket\mathcal{F}\rrbracket$ that approximate a classical function with probability at least $p$.\par
 
% END S4.SS2.p2
 
% BEGIN Thmtheorem11
 
\par\medskip\noindent\textbf{Theorem 11 (Soundness and Completeness).}\par
 
% BEGIN Thmtheorem11.p1
 
$\llbracket\textnormal{{pbp}}{}\rrbracket_{\geq\frac{2}{3}}=\textnormal{{fbqp}}{}$.\par
 
% END Thmtheorem11.p1
 
% END Thmtheorem11
 
% END S4.SS2
 
% BEGIN S4.SS3
 
\par\medskip\noindent\textbf{4.3 Extension to Programs in $\textnormal{{wf}}\cap\textnormal{{width}}_{\leq 1}$}\par
 
% BEGIN S4.SS3.p1
 
The basic restriction of pbp can be thought of as ensuring that qubits are consumed in a uniform way. As a consequence, any two procedure calls on inputs of size $n$ will correspond to precisely the same $n$ qubits. This guarantees that two compatible procedure calls can be merged simply by combining control structures in the same ancilla.\par
 
% END S4.SS3.p1
 
% BEGIN S4.SS3.p2
 
Without this constraint (i.e., on $\textnormal{{wf}}\cap\textnormal{{width}}_{\leq 1}$), we can use precisely the same compilation strategy but where merging is done by control-swapping registers into a single procedure call. For $k$ instances of a procedure on input size $n$, this requires $k-1$ controlled swaps, each requiring $O(n)$ operations in the worst case. An example of merging two unitary gates $U$ with controlled-swaps is given in Figure 12. If we allow for parallelization of gates in the circuit, this can be done in logarithmic depth, as shown in Lemma 12. An example of a controlled permutation is shown in Figure 13, where vertical dashed lines separate time slices where gates can be applied simultaneously.\par
 
% END S4.SS3.p2
 
% BEGIN S4.F12

\begin{center}\includegraphics[width=0.98\linewidth,height=0.70\textheight,keepaspectratio]{graphic-18.pdf}\end{center}
 
Figure 12: Merging orthogonal controlled statements.\par
 
% END S4.F12
 
% BEGIN S4.F13

\begin{center}\includegraphics[width=0.98\linewidth,height=0.70\textheight,keepaspectratio]{graphic-19.pdf}\end{center}
 
Figure 13: Implementation of a controlled permutation in log-depth.\par
 
% END S4.F13
 
% BEGIN Thmtheorem12
 
\par\medskip\noindent\textbf{Lemma 12.}\par
 
% BEGIN Thmtheorem12.p1
 
Any controlled permutation on $n$ qubits can be performed with a quantum circuit of size $O(n)$ and depth $O(\log n)$.\par
 
% END Thmtheorem12.p1
 
% END Thmtheorem12
 
% BEGIN S4.SS3.p3
 
An extension of compile with controlled-swapping as in Figure 12 used for merging procedures ensures the following bound for $\textnormal{{wf}}\cap\textnormal{{width}}_{\leq 1}$ programs.\par
 
% END S4.SS3.p3
 
% BEGIN Thmtheorem13
 
\par\medskip\noindent\textbf{Theorem 13.}\par
 
% BEGIN Thmtheorem13.p1
 
For $\mathrm{P}\in\textnormal{{wf}}\cap\textnormal{{width}}_{\leq 1}$, $\textnormal{{{compile}}{}}(\mathrm{P},n)$ results in a quantum circuit of size $O(n\cdot\textnormal{Time}_{\mathrm{P}}(n))$ and depth ${O}(\log(n)\cdot\textnormal{Time}_{\mathrm{P}}(n))$.\par
 
% END Thmtheorem13.p1
 
% END Thmtheorem13
 
% END S4.SS3
 
% BEGIN S4.SS4
 
\par\medskip\noindent\textbf{4.4 Comparison with Existing Work}\par
 
% BEGIN S4.SS4.p1
 
The compile algorithm strictly improves upon the size bounds of other compilation algorithms [17], as illustrated by Table 1. These bounds can be empirically tested using the implementation of the compiler available at gitlab.inria.fr/mmachado/pfoq-compiler. In some cases, such as Examples 16 and 15 given in Appendix B, we find families of programs where the gap in complexity (i.e. the degree of the polynomial) can be made arbitrarily large.\par
 
% END S4.SS4.p1
 
% BEGIN S4.SS4.p2
 
In $k$-Chained Substrings (Example 16), we consider the problem of detecting whether an input string contains substrings $y_{1},\dots,y_{k}$ appearing in that order. For a certain choice of substrings $y_{i}$, we define a pbp program of linear time complexity for which the compilation strategy in [17] results in circuits of size $\Theta(n^{2k})$. The problem Sum($r$), given in Example 17, consists of checking if an input $x=x_{1}\dots x_{n}$ satisfies $\sum_{i=1}^{n}x_{i}=r$.\par
 
% END S4.SS4.p2
 
% BEGIN S4.T1
 
Table 1: Circuit size complexity bounds.\par
     Circuit complexity     Problem   Example     [17]     compile     Full Adder   Example 15     $\Theta(n)$     $\Theta(n)$     Quantum Fourier Transform   Example 6     $\Theta(n^{2})$     $\Theta(n^{2})$      $k$-Chained Substrings   Example 16     $\Theta(n^{2k})$     $\Theta(n)$     Sum($r$), $r\geq 1$    Example 17     $\Theta(n^{r})$     $\Theta(n)$      
% END S4.T1
 
% END S4.SS4
 
% END S4


% BEGIN A1
 
\par\medskip\noindent\textbf{Appendix A Omitted proofs}\par
 
% BEGIN A1.p1
 
% BEGIN Thmtheorem8
 
\par\medskip\noindent\textbf{Theorem 8 (Soundness of compilation). [Restated, see original statement.]}\par
 
% BEGIN Thmtheorem8.p1
 
Given a program $\mathrm{P}\in\textnormal{{wf}}\cap\textnormal{{width}}_{\leq 1}$ and a quantum state $\left|\psi\right\rangle\in\mathcal{H}_{2^{n}}$ we have that $\llbracket\textnormal{{compile}}{}(\mathrm{P},n)\rrbracket(\left|\psi\right%
\rangle)=\llbracket\mathrm{P}\rrbracket(\left|\psi\right\rangle)$.\par
 
% END Thmtheorem8.p1
 
% END Thmtheorem8
 
% END A1.p1
 
% BEGIN A1.SS0.SSS0.P0.SPx1
 
\par\medskip\noindent\textbf{Proof.}\par
 
% BEGIN A1.SS0.SSS0.P0.SPx1.p1
 
By induction on the structure of the program $\mathrm{P}=\mathrm{D}::\mathrm{S}$ in $\textnormal{{wf}}\cap\textnormal{{width}}_{\leq 1}$. One can check by inspection of each case that the compile rules for non-recursive statements corresponds to the straightforward circuit semantics of quantum programs.\par
 
% END A1.SS0.SSS0.P0.SPx1.p1
 
% BEGIN A1.SS0.SSS0.P0.SPx1.p2
 
Likewise, the rewriting rules of optimize, given in Figure 9, can be easily checked using the orthogonality invariant of Lemma 7. For instance, consider the case of the sequential statement $\mathrm{S}=\mathrm{S}_{1}\ \mathrm{S}_{2}$. In the case where $w^{\text{\color[rgb]{0.79296875,0.3203125,0.8359375}\definecolor[named]{%
pgfstrokecolor}{rgb}{0.79296875,0.3203125,0.8359375}{{proc}}}}_{\mathrm{P}}(%
\mathrm{S}_{1})=1$, we can derive the rule for the statement with the following steps. We omit the label $l_{\mathrm{Cst}}$ to simplify the reading.\par
 
% END A1.SS0.SSS0.P0.SPx1.p2
 
% BEGIN A1.SS0.SSS0.P0.SPx1.p3
 
\begin{center}\includegraphics[width=0.98\linewidth,height=0.70\textheight,keepaspectratio]{graphic-20.pdf}\end{center}
 
% END A1.SS0.SSS0.P0.SPx1.p3
 
% BEGIN A1.SS0.SSS0.P0.SPx1.p4
 
where $(a)$ corresponds to the definition of the sequential statement, in $(b)$ we make use of the fact that $(cs,\mathrm{S}_{2})$ is orthogonal to all controlled statements in $l_{\mathrm{Cst}}$ and therefore can be commuted in the circuit, and in step $(c)$ we likewise make use of orthogonality, in this case with the controlled statements in the contextual list. Other case can be inspected to follow a short sequence of steps, such as described for the sequential statement.\par
 
% END A1.SS0.SSS0.P0.SPx1.p4
 
% BEGIN A1.SS0.SSS0.P0.SPx1.p5
 
For instance, in the case of anchoring, we consider the following composition of rules\par
 
% END A1.SS0.SSS0.P0.SPx1.p5
 
% BEGIN A1.SS0.SSS0.P0.SPx1.p6
 
\begin{center}\includegraphics[width=0.98\linewidth,height=0.70\textheight,keepaspectratio]{graphic-21.pdf}\end{center}
 
% END A1.SS0.SSS0.P0.SPx1.p6
 
% BEGIN A1.SS0.SSS0.P0.SPx1.p7
 
where, likewise, $(a)$ corresponds to the typical circuit semantics of a procedure call (with an added anchoring ancilla), and $(b)$ and $(c)$ make use of the orthogonality of $cs$ with the right side of the circuit. The validity of merging and also be checked:\par
 
% END A1.SS0.SSS0.P0.SPx1.p7
 
% BEGIN A1.SS0.SSS0.P0.SPx1.p8
 
\begin{center}\includegraphics[width=0.98\linewidth,height=0.70\textheight,keepaspectratio]{graphic-22.pdf}\end{center}
 
% END A1.SS0.SSS0.P0.SPx1.p8
 
% BEGIN A1.SS0.SSS0.P0.SPx1.p9
 
Step $(a)$ can be seen as the definition of the procedure call. Since we are in the merging scenario, a similar procedure call has already been performed and is anchored to its corresponding ancilla. Without loss of generality (since all controlled statements in $l_{\mathrm{Cst}}$ are orthogonal and therefore commute) we assume that this call appears at the end of $l_{\mathrm{Cst}}$. Rule $(b)$ simply indicates that since $cs$ is orthogonal to other control structures in $l_{\mathrm{Cst}}$, we may move the controlled statements so that they are adjacent, and where we apply step $(c)$ to perform merging. Since $cs$ is orthogonal to $\cdot[a_{|\mathrm{s}|}^{\text{\color[rgb]{0.79296875,0.3203125,0.8359375}%
\definecolor[named]{pgfstrokecolor}{rgb}{0.79296875,0.3203125,0.8359375}{{proc%
'}}}}=1]$ the control is added to the ancilla as expected. Finally, orthogonality of $cs$ with other control structures means that we may move the two controlled-$NOT$s to the edges of the circuit.\par
 
% END A1.SS0.SSS0.P0.SPx1.p9
 
% BEGIN A1.SS0.SSS0.P0.SPx1.p10
 
We now consider the validity of the rule for the contextual circuit. Consider a list of mutually-orthogonal controlled statements, where the sequential form of $(cs_{i},\mathrm{S}_{i})$ is given by the lists of controlled statements $l^{(i)}_{1},\dots,l^{(i)}_{t}$, then we have that:\par
 
% END A1.SS0.SSS0.P0.SPx1.p10
 
% BEGIN A1.SS0.SSS0.P0.SPx1.p11
 
\begin{center}\includegraphics[width=0.98\linewidth,height=0.70\textheight,keepaspectratio]{graphic-23.pdf}\end{center}
 
% END A1.SS0.SSS0.P0.SPx1.p11
 
% BEGIN A1.SS0.SSS0.P0.SPx1.p12
 
The final step comes from the fact that all $l^{(1)}_{j}$ are orthogonal to $cs_{2}$, …, $cs_{k}$ since they include $cs_{1}$. This can be done for all other $(cs_{i},\mathrm{S}_{i})$, and in the end we obtain the arrangement $l_{1}^{(1)}\ldots l_{1}^{(k)}l_{2}^{(1)}\ldots l_{2}^{(k)}\ldots l_{t}^{(1)}%
\ldots l_{t}^{(k)}$.\par
 
% END A1.SS0.SSS0.P0.SPx1.p12
 
% BEGIN A1.SS0.SSS0.P0.SPx1.p13
 
Given a certain $1\leq j\leq t$, we have that all controlled statements in $\cup_{i=1}^{k}l_{j}^{(i)}$ (that is, all controlled statements occurring in time $j$) are pairwise orthogonal. Therefore, we may rearrange their order according to their recursivity, and in doing so we may consider each time separately. For instance, in time $1$, let $\mathcal{L}\triangleq\cup_{i=1}^{k}l_{1}^{(i)}$, and let $\text{\color[rgb]{0.79296875,0.3203125,0.8359375}\definecolor[named]{%
pgfstrokecolor}{rgb}{0.79296875,0.3203125,0.8359375}{{proc${}_{1}$}}},\,\text{%
\color[rgb]{0.79296875,0.3203125,0.8359375}\definecolor[named]{pgfstrokecolor}%
{rgb}{0.79296875,0.3203125,0.8359375}{{proc${}_{2}$}}},\dots\text{\color[rgb]{%
0.79296875,0.3203125,0.8359375}\definecolor[named]{pgfstrokecolor}{rgb}{%
0.79296875,0.3203125,0.8359375}{{proc${}_{m}$}}}$ denote procedures belonging to different recursion families. Then, we perform the partition:\par
    $\displaystyle\mathcal{L}_{0}$ $\displaystyle\triangleq\{(cs,\mathrm{S})\in\mathcal{L}:\not\exists\text{\color%
[rgb]{0.79296875,0.3203125,0.8359375}\definecolor[named]{pgfstrokecolor}{rgb}{%
0.79296875,0.3203125,0.8359375}{{proc}}}\text{ such that }w_{\mathrm{P}}^{%
\text{\color[rgb]{0.79296875,0.3203125,0.8359375}\definecolor[named]{%
pgfstrokecolor}{rgb}{0.79296875,0.3203125,0.8359375}{{proc}}}}(\mathrm{S}^{%
\text{\color[rgb]{0.79296875,0.3203125,0.8359375}\definecolor[named]{%
pgfstrokecolor}{rgb}{0.79296875,0.3203125,0.8359375}{{proc}}}})=1\text{ and }w%
_{\mathrm{P}}^{\text{\color[rgb]{0.79296875,0.3203125,0.8359375}\definecolor[%
named]{pgfstrokecolor}{rgb}{0.79296875,0.3203125,0.8359375}{{proc}}}}(\mathrm{%
S})=1\}.$     $\displaystyle\mathcal{L}_{i}$ $\displaystyle\triangleq\{(cs,\mathrm{S})\in\mathcal{L}:w_{\mathrm{P}}^{\text{%
\color[rgb]{0.79296875,0.3203125,0.8359375}\definecolor[named]{pgfstrokecolor}%
{rgb}{0.79296875,0.3203125,0.8359375}{{proc${}_{i}$}}}}(\mathrm{S})=1\},\quad i%
=1,\dots,m.$    
By the definition of the sequential form of each controlled statement, we note that the partition is well-defined (e.g. there are no statements containing calls to more than one procedure). Therefore, we are able to rearrange $\mathcal{L}$ and perform calls to optimize on each $\mathcal{L}_{i\geq 1}$ and a call to compile on $\mathcal{L}_{0}$. After composing the obtained circuits, we get the rule given in Section 3. This concludes the proof. $\hfill\blacktriangleleft$\par
 
% END A1.SS0.SSS0.P0.SPx1.p13
 
% BEGIN A1.SS0.SSS0.P0.SPx1.p14
 
We introduce the notion of rank of a procedure for use in the proof of Theorem 9.\par
 
% END A1.SS0.SSS0.P0.SPx1.p14
 
% BEGIN Thmtheorem14
 
\par\medskip\noindent\textbf{Definition 14 (Rank).}\par
 
% BEGIN Thmtheorem14.p1
 
Set $\max(\emptyset)\triangleq 0$. Given a program $\mathrm{P}$, the rank of a procedure $\text{\color[rgb]{0.79296875,0.3203125,0.8359375}\definecolor[named]{%
pgfstrokecolor}{rgb}{0.79296875,0.3203125,0.8359375}{{proc}}}\in\mathrm{P}{}$, denoted $rk_{\mathrm{P}}(\text{\color[rgb]{0.79296875,0.3203125,0.8359375}\definecolor[%
named]{pgfstrokecolor}{rgb}{0.79296875,0.3203125,0.8359375}{{proc}}})$, is defined as follows:\par
 
% END Thmtheorem14.p1
 
% BEGIN Thmtheorem14.p2
    \[rk_{\mathrm{P}}(\text{\color[rgb]{0.79296875,0.3203125,0.8359375}\definecolor[%
named]{pgfstrokecolor}{rgb}{0.79296875,0.3203125,0.8359375}{{proc}}})%
\triangleq\begin{cases}0,&\text{ if }\neg(\exists\text{\color[rgb]{%
0.79296875,0.3203125,0.8359375}\definecolor[named]{pgfstrokecolor}{rgb}{%
0.79296875,0.3203125,0.8359375}{{proc'}}},\ \text{\color[rgb]{%
0.79296875,0.3203125,0.8359375}\definecolor[named]{pgfstrokecolor}{rgb}{%
0.79296875,0.3203125,0.8359375}{{proc}}}\succeq_{\mathrm{P}}\text{\color[rgb]{%
0.79296875,0.3203125,0.8359375}\definecolor[named]{pgfstrokecolor}{rgb}{%
0.79296875,0.3203125,0.8359375}{{proc'}}}),\\
\max_{\text{\color[rgb]{0.79296875,0.3203125,0.8359375}\definecolor[named]{%
pgfstrokecolor}{rgb}{0.79296875,0.3203125,0.8359375}{{proc}}}\succeq_{\mathrm{%
P}}\text{\color[rgb]{0.79296875,0.3203125,0.8359375}\definecolor[named]{%
pgfstrokecolor}{rgb}{0.79296875,0.3203125,0.8359375}{{proc'}}}}\{rk_{\mathrm{P%
}}(\text{\color[rgb]{0.79296875,0.3203125,0.8359375}\definecolor[named]{%
pgfstrokecolor}{rgb}{0.79296875,0.3203125,0.8359375}{{proc'}}})\},&\text{ if }%
\exists\text{\color[rgb]{0.79296875,0.3203125,0.8359375}\definecolor[named]{%
pgfstrokecolor}{rgb}{0.79296875,0.3203125,0.8359375}{{proc'}}},\ \text{\color[%
rgb]{0.79296875,0.3203125,0.8359375}\definecolor[named]{pgfstrokecolor}{rgb}{%
0.79296875,0.3203125,0.8359375}{{proc}}}\succeq_{\mathrm{P}}\text{\color[rgb]{%
0.79296875,0.3203125,0.8359375}\definecolor[named]{pgfstrokecolor}{rgb}{%
0.79296875,0.3203125,0.8359375}{{proc'}}}\wedge\neg(\text{\color[rgb]{%
0.79296875,0.3203125,0.8359375}\definecolor[named]{pgfstrokecolor}{rgb}{%
0.79296875,0.3203125,0.8359375}{{proc}}}\sim_{\mathrm{P}}\text{\color[rgb]{%
0.79296875,0.3203125,0.8359375}\definecolor[named]{pgfstrokecolor}{rgb}{%
0.79296875,0.3203125,0.8359375}{{proc}}}),\\
1+\max_{\text{\color[rgb]{0.79296875,0.3203125,0.8359375}\definecolor[named]{%
pgfstrokecolor}{rgb}{0.79296875,0.3203125,0.8359375}{{proc}}}\succ_{\mathrm{P}%
}\text{\color[rgb]{0.79296875,0.3203125,0.8359375}\definecolor[named]{%
pgfstrokecolor}{rgb}{0.79296875,0.3203125,0.8359375}{{proc'}}}}\{rk_{\mathrm{P%
}}(\text{\color[rgb]{0.79296875,0.3203125,0.8359375}\definecolor[named]{%
pgfstrokecolor}{rgb}{0.79296875,0.3203125,0.8359375}{{proc'}}})\},&\text{ if }%
\text{\color[rgb]{0.79296875,0.3203125,0.8359375}\definecolor[named]{%
pgfstrokecolor}{rgb}{0.79296875,0.3203125,0.8359375}{{proc}}}\sim_{\mathrm{P}}%
\text{\color[rgb]{0.79296875,0.3203125,0.8359375}\definecolor[named]{%
pgfstrokecolor}{rgb}{0.79296875,0.3203125,0.8359375}{{proc}}}.\end{cases}\]    
% END Thmtheorem14.p2
 
% END Thmtheorem14
 
% BEGIN A1.SS0.SSS0.P0.SPx1.p15
 
% BEGIN Thmtheorem9
 
\par\medskip\noindent\textbf{Theorem 9 (No branch sequentialization). [Restated, see original statement.]}\par
 
% BEGIN Thmtheorem9.p1
 
Given a program $\mathrm{P}\in\textnormal{{pbp}}{}$ and input $n\in\mathbb{N}$, we have that $\#\textnormal{{{compile}}{}}(\mathrm{P},n)=O(\textnormal{Time}_{\mathrm{P}}(n))$ holds.\par
 
% END Thmtheorem9.p1
 
% END Thmtheorem9
 
% END A1.SS0.SSS0.P0.SPx1.p15
 
% END A1.SS0.SSS0.P0.SPx1
 
% BEGIN A1.SS0.SSS0.P0.SPx2
 
\par\medskip\noindent\textbf{Proof.}\par
 
% BEGIN A1.SS0.SSS0.P0.SPx2.p1
 
The theorem can be shown by structural induction on the program body. All cases are straightforward except the one of the quantum control case. We make use of the following two facts regarding pbp programs:\par
  
\par\noindent  (a) 
% BEGIN A1.I1.i1.p1
 
The size of the circuit is directly proportional to its total number of unique procedure calls (in the sense required for merging), and\par
 
% END A1.I1.i1.p1
 
 
\par\noindent  (b) 
% BEGIN A1.I1.i2.p1
 
a recursive procedure call results in $O(n)$ unique calls to procedures of the same rank. This is because unique calls may only differ on procedure name (of which there is a constant amount) or input size (for which there is a linear number of possibilites).\par
 
% END A1.I1.i2.p1
 
  
% END A1.SS0.SSS0.P0.SPx2.p1
 
% BEGIN A1.SS0.SSS0.P0.SPx2.p2
 
Consider a quantum control statement $\mathrm{S}=\text{\color[rgb]{0.03515625,0.0390625,0.72265625}\definecolor[%
named]{pgfstrokecolor}{rgb}{0.03515625,0.0390625,0.72265625}{{qcase }}}\mathrm%
{q}\text{\color[rgb]{0.03515625,0.0390625,0.72265625}\definecolor[named]{%
pgfstrokecolor}{rgb}{0.03515625,0.0390625,0.72265625}{{ of }}}\{0\to\mathrm{S}%
_{0},\,1\to\mathrm{S}_{1}\}$ appearing in optimize in the context of a recursive procedure proc. By (a), the circuit size for $\mathrm{S}_{0}$ and $\mathrm{S}_{1}$ are proportional to the (total) number of procedure calls in each statement, separately. We check that the number of ancillas created for $\mathrm{S}$ is bounded by the maximum number of ancillas between $\mathrm{S}_{0}$ and $\mathrm{S}_{1}$.\par
 
% END A1.SS0.SSS0.P0.SPx2.p2
 
% BEGIN A1.SS0.SSS0.P0.SPx2.p3
 
We proceed by induction on the rank $r$ of the procedure. The base case of $r=1$ is given by (b), and so we may consider $r>1$. For the inductive case, we consider two scenarios:\par
  
\par\noindent  $\blacksquare$ 
% BEGIN A1.I2.i1.p1
 
$w_{\text{\color[rgb]{0.79296875,0.3203125,0.8359375}\definecolor[named]{%
pgfstrokecolor}{rgb}{0.79296875,0.3203125,0.8359375}{{proc}}}}^{\mathrm{P}}(%
\mathrm{S}_{0})=w_{\text{\color[rgb]{0.79296875,0.3203125,0.8359375}%
\definecolor[named]{pgfstrokecolor}{rgb}{0.79296875,0.3203125,0.8359375}{{proc%
}}}}^{\mathrm{P}}(\mathrm{S}_{1})=1$. In this case, both $\mathrm{S}_{0}$ and $\mathrm{S}_{1}$ are of rank $r$, and all their recursive procedure calls may be merged. Therefore, the asymptotic number of such calls is bounded between the maximum between $\mathrm{S}_{0}$ and $\mathrm{S}_{1}$ (consider that, if there is no overlap between the ancillas needed, their number is still bounded linearly). Applying the IH on the procedures of rank $r-1$ we obtain the desired result.\par
 
% END A1.I2.i1.p1
 
 
\par\noindent  $\blacksquare$ 
% BEGIN A1.I2.i2.p1
 
$w_{\text{\color[rgb]{0.79296875,0.3203125,0.8359375}\definecolor[named]{%
pgfstrokecolor}{rgb}{0.79296875,0.3203125,0.8359375}{{proc}}}}^{\mathrm{P}}(%
\mathrm{S}_{0})=0$ and $w_{\text{\color[rgb]{0.79296875,0.3203125,0.8359375}\definecolor[named]{%
pgfstrokecolor}{rgb}{0.79296875,0.3203125,0.8359375}{{proc}}}}^{\mathrm{P}}(%
\mathrm{S}_{1})=1$. In this case, $\mathrm{S}_{0}$ contains calls to procedures of rank $r^{\prime}<r$ whereas $\mathrm{S}_{1}$ contains calls to procedures of rank $r$. According to optimize, statement $\mathrm{S}_{0}$ is kept in the contextual circuit until no more statements are recursive relative to proc. The statements of rank $r^{\prime}$ which are present in $\mathrm{S}_{0}$ are then merged with the equivalent procedures that were derived from $\mathrm{S}_{1}$ and also added to $l_{\text{M}}$. The number of procedures of rank $r^{\prime}$ is bounded asymptotically by the maximum between those in $\mathrm{S}_{0}$ and $\mathrm{S}_{1}$, therefore we obtain our result.\par
 
% END A1.I2.i2.p1
 
  
$\hfill\blacktriangleleft$\par
 
% END A1.SS0.SSS0.P0.SPx2.p3
 
% BEGIN A1.SS0.SSS0.P0.SPx2.p4
 
% BEGIN Thmtheorem11
 
\par\medskip\noindent\textbf{Theorem 11 (Soundness and Completeness). [Restated, see original statement.]}\par
 
% BEGIN Thmtheorem11.p1
 
$\llbracket\textnormal{{pbp}}{}\rrbracket_{\geq\frac{2}{3}}=\textnormal{{fbqp}}{}$.\par
 
% END Thmtheorem11.p1
 
% END Thmtheorem11
 
% END A1.SS0.SSS0.P0.SPx2.p4
 
% END A1.SS0.SSS0.P0.SPx2
 
% BEGIN A1.SS0.SSS0.P0.SPx3
 
\par\medskip\noindent\textbf{Proof.}\par
 
% BEGIN A1.SS0.SSS0.P0.SPx3.p1
 
Since $\textnormal{{pbp}}\subsetneq\textnormal{{wf}}\cap\textnormal{{width}}_{\leq 1}%
\subsetneq\textnormal{{pfoq}}{}$, with pfoq the language of [17], we have that $\llbracket\textnormal{{pbp}}{}\rrbracket\subseteq\llbracket\textnormal{{wf}}%
\cap\textnormal{{width}}_{\leq 1}\rrbracket\subseteq\llbracket\textnormal{{%
pfoq}}{}\rrbracket$ and, by pfoq soundness [17, Theorem 3], it also holds that $\llbracket\textnormal{{pbp}}{}\rrbracket_{\geq\frac{2}{3}}\subseteq\textnormal%
{{fbqp}}{}$. Completeness can be proven by showing that pbp can simulate the function algebra in [29], known to be complete for quantum polynomial time. The proof can be done using the same construction as in [17, Theorem 5].$\hfill\blacktriangleleft$\par
 
% END A1.SS0.SSS0.P0.SPx3.p1
 
% BEGIN A1.SS0.SSS0.P0.SPx3.p2
 
See 12\par
 
% END A1.SS0.SSS0.P0.SPx3.p2
 
% END A1.SS0.SSS0.P0.SPx3
 
% BEGIN A1.SS0.SSS0.P0.SPx4
 
\par\medskip\noindent\textbf{Proof.}\par
 
% BEGIN A1.SS0.SSS0.P0.SPx4.p1
 
Any permutation can be written as the composition of two sets of disjoint transpositions, and therefore any permutation can be performed in constant time, using two time steps [24]. To perform a controlled permutation, it suffices to create $O(n)$ ancillas with the correct controlled state, which can be done in $O(\log n)$ depth with $O(n)$ gates. $\hfill\blacktriangleleft$\par
 
% END A1.SS0.SSS0.P0.SPx4.p1
 
% END A1.SS0.SSS0.P0.SPx4
 
% END A1
